%  LaTeX support: latex@mdpi.com 
%  For support, please attach all files needed for compiling as well as the log file, and specify your operating system, LaTeX version, and LaTeX editor.

%=================================================================
\documentclass[make,article,submit,moreauthors]{Definitions/mdpi} 
%\documentclass[journal,article,submit,moreauthors,unicode]{Definitions/mdpi} % For some letters Palatino lacks
%\documentclass[preprints,article,submit,moreauthors]{Definitions/mdpi} 
%\documentclass[preprints,article,submit,moreauthors,unicode]{Definitions/mdpi} 
% For posting an early version of this manuscript as a preprint, you may use "preprints" as the journal. Changing "submit" to "accept" before posting will remove line numbers.

\setlength{\headheight}{18.2pt}
\addtolength{\topmargin}{-6.2pt}

%--------------------
% Class Options:
%--------------------
%----------
% journal
%----------
% Choose between the following MDPI journals:
% accountaudit, acn, acoustics, actuators, addictions, addictprev, adhesives, adolescents, aee, aerobiology, aeronautics, aerospace, agriculture, agriengineering, agrochemicals, agronomy, ai, aiagents, aichem, aieng, aimater, aimed, aipa, air, aisens, aisoc, algorithms, allergies, alloys, amh, analog, analytica, analytics, anatomia, anesthres, animals, antibiotics, antibodies, antioxidants, applbiosci, appliedchem, appliedmath, appliedphys, applmech, applmicrobiol, applnano, applsci, aps, aquacj, archaeolstud, architecture, arm, arthropoda, asc, asd, asi, astronautics, astronomy, atmosphere, atoms, audiolres, automation, axioms, bacteria, batteries, bcrc, bdcc, beverages, bioactives, biochem, bioengineering, biologics, biology, biomass, biomechanics, biomed, biomedicines, biomedinformatics, biomimetics, biomolecules, biophysica, bioresourbioprod, biosensors, biosphere, biotech, birds, blockchains, bloods, blsf, brainsci, breath, buildings, cancers, carbon, cardio, cardiogenetics, cardiovascmed, catalysts, cells, ceramics, challenges, chemeco, chemengineering, chemistry, chemosensors, chemproc, children, chips, cimb, civileng, cleantechnol, climate, clinbioenerg, clinpract, clockssleep, cmd, cmtr, coasts, coatings, colloids, colorants, commodities, complexities, complications, compounds, computation, computers, condensedmatter, conservation, constrmater, cosmetics, covid, crops, cryo, cryptography, crystals, csmf, ctn, culture, curroncol, cyber, dairy, data, ddc, dentistry, dermato, dermatopathology, designs, devices, dgg, dhi, diabetology, diagnostics, dietetics, digital, disabilities, diseases, diversity, dna, drones, drugreact, drylands, dynamics, earth, ebj, ece, ecm, ecologies, edm, eesp, electricity, electrochem, electronicmat, electronics, embodiedintell, encyclopedia, endocrines, energies, eng, engproc, entomology, entropy, environments, environremediat, enzymology, epidemiologia, epigenomes, epilepsy, esa, est, fermentation, fibers, fintech, fire, fishes, fluids, foodeng, foods, forecasting, forensicsci, forests, fossstud, foundations, fractalfract, freshwater, fuels, future, futureinternet, futurepharmacol, futurephys, futuretransp, galaxies, gases, gastroent, gastrointestdisord, gastronomy, gels, genes, geochem, geographies, geohazards, geomatics, geometry, geosciences, geotechnics, geriatrics, germs, glacies, grainsci, grasses, green, greenhealth, gucdd, gynecolcancer, hardware, hazardousmatters, healthcare, healthcmater, hearts, hemato, hematolrep, hep, heritage, higheredu, highthroughput, horticulturae, hospitals, humanoids, hydrobiology, hydrogen, hydrology, hydrometeorology, hydropower, hygiene, idr, iic, ijem, ijerph, ijgi, ijmd, ijms, ijns, ijom, ijp, ijpb, ijt, ijtm, ijtpp, ijtst, ime, immuno, industries, inflammj, informatics, information, infrastructures, inorganics, insects, instruments, inventions, iot, iron, j, jaestheticmed, jal, japma, jcdd, jcm, jcp, jcrm, jcs, jcto, jdad, jdb, jdcp, jdream, jemr, jeta, jfb, jfmk, jfth, jgbg, jgg, jimaging, jlpea, jmahp, jmmp, jmms, jmp, jmse, jne, jnt, jof, joitmc, joma, jop, joptm, jor, jox, jpbi, jphytomed, jpm, jsan, jtaer, jvd, jzbg, kidneydial, kinasesphosphatases, knowledge, labmed, laboratories, lam, land, langtechnol, lasers, legumes, life, lights, limnolrev, lipidology, liquids, livers, logics, logistics, lubricants, lymphatics, machines, macromol, magnetism, magnetochemistry, make, marinedrugs, maritime, materials, materproc, mathematics, mca, mcb, measurements, meb, medicina, medicines, medsci, membranes, merits, metabolites, metals, metamaterialsadv, meteorology, methane, metrics, metrology, micro, microarrays, microbiolres, microelectronics, micromachines, microorganisms, microplastics, microwave, minerals, mining, mmphys, modelling, molbank, molecules, mollusca, mps, msf, mti, multimedia, muscles, mutation, nanoenergyadv, nanomanufacturing, nanomaterials, ncrna, ndt, network, neuroglia, neuroimaging, neurolint, neurosci, nitrogen, notspecified, nursrep, nutraceuticals, nutrients, obesities, occuphealth, oceans, ohbm, onco, optics, optoelectronics, oral, organics, organoids, osteology, oxygen, pandemics, parasites, parasitologia, particles, pathogens, pathophysiology, pediatrrep, pets, pharmaceuticals, pharmaceutics, pharmacoepidemiology, pharmacy, phc, philosophies, photochem, photonics, photovoltaics, phycology, physchem, physics, physiologia, plants, plasma, platforms, pnr, pollutants, polymers, polysaccharides, populations, poultry, powders, precision, precisoncol, preprints, prevhealthc, proceedings, processes, prosthesis, proteomes, psf, psychiatryint, psychoactives, purification, quantumrep, quaternary, qubs, radiation, rareearths, rdt, reactions, realestate, receptors, recycling, regeneration, remotesensing, reports, reprodmed, resources, rheumato, rjpm, robotics, rsee, ruminants, safety, sanpp, sca, sci, scipharm, sclerosis, seeds, sensors, senssyst, separations, sexes, shi, signals, sinusitis, siuj, skins, smartcities, smartfish, smartgrids, sna, societies, software, soilsystems, solar, solids, spectroscj, sports, standards, stats, std, stratsediment, stresses, strokerep, superintelligence, surfaces, surgeries, susagric, suschem, sustainability, sustainfoods, sustainpackag, symmetry, synbio, systems, tae, targets, taxonomy, technologies, telecom, test, textiles, thalassrep, therapeutics, thermo, timespace, tomography, toxics, toxins, tph, transplantology, transportation, traumacare, traumas, tri, tropicalmed, universe, urbansci, uro, vaccines, vehicles, venereology, vetsci, vibration, virtualworlds, viruses, vision, waste, water, welding, wem, wevj, wild, wind, women, world, zoonoticdis

%---------
% article
%---------
% The default type of manuscript is "article", but can be replaced by: 
% abstract, addendum, article, benchmark, book, bookreview, briefcommunication, briefreport, casereport, changes, clinicopathologicalchallenge, comment, commentary, communication, conceptpaper, conferenceproceedings, correction, conferencereport, creative, datadescriptor, discussion, entry, expressionofconcern, extendedabstract, editorial, essay, erratum, fieldguide, hypothesis, interestingimages, letter, meetingreport, monograph, newbookreceived, obituary, opinion, proceedingpaper, projectreport, reply, retraction, review, perspective, protocol, shortnote, studyprotocol, supfile, systematicreview, technicalnote, viewpoint, guidelines, registeredreport, tutorial, giantsinurology, urologyaroundtheworld, patentsummary

% supfile = supplementary materials

%----------
% submit
%----------
% The class option "submit" will be changed to "accept" by the Editorial Office when the paper is accepted. This will only make changes to the frontpage (e.g., the logo of the journal will get visible), the headings, and the copyright information. Also, line numbering will be removed. Journal info and pagination for accepted papers will also be assigned by the Editorial Office.

%------------------
% moreauthors
%------------------
% If there is only one author the class option oneauthor should be used. Otherwise use the class option moreauthors.

%=================================================================
% MDPI internal commands - do not modify
\firstpage{1} 
\makeatletter 
\setcounter{page}{\@firstpage} 
\makeatother
\pubvolume{1}
\issuenum{1}
\articlenumber{0}
\pubyear{2026}
\copyrightyear{2026}
%\externaleditor{Firstname Lastname} % More than 1 editor, please add `` and '' before the last editor name
\datereceived{ } 
\daterevised{ } % Comment out if no revised date
\dateaccepted{ } 
\datepublished{ } 
%\datecorrected{} % For corrected papers: "Corrected: XXX" date in the original paper.
%\dateretracted{} % For retracted papers: "Retracted: XXX" date in the original paper.
%\doinum{} % Used for some special journals, like molbank
%\pdfoutput=1 % Uncommented for upload to arXiv.org
%\CorrStatement{yes}  % For updates
%\longauthorlist{yes} % For many authors that exceed the left citation part
%\IsAssociation{yes} % For association journals

%=================================================================
% Add packages and commands here. The following packages are loaded in our class file: fontenc, inputenc, calc, indentfirst, fancyhdr, graphicx, epstopdf, lastpage, ifthen, float, amsmath, amssymb, lineno, setspace, enumitem, mathpazo, booktabs, titlesec, etoolbox, tabto, xcolor, colortbl, soul, multirow, microtype, tikz, totcount, changepage, attrib, upgreek, array, tabularx, pbox, ragged2e, tocloft, marginnote, marginfix, enotez, amsthm, natbib, hyperref, cleveref, scrextend, url, geometry, newfloat, caption, draftwatermark, seqsplit
% cleveref: load \crefname definitions after \begin{document}

\usepackage{subfig}
\usepackage{longtable}

%=================================================================
% Please use the following mathematics environments: Theorem, Lemma, Corollary, Proposition, Characterization, Property, Problem, Example, ExamplesandDefinitions, Hypothesis, Remark, Definition, Notation, Assumption
%% For proofs, please use the proof environment (the amsthm package is loaded by the MDPI class).

%=================================================================
% Full title of the paper (Capitalized)
\Title{Zero-Inflated Time Series Generation}

% Author OrchiZero-Inflated Time Series Generationd ID: enter ID or remove command
\newcommand{\orcidauthorA}{0000-0002-0098-4228} % Add \orcidA{} behind the author's name
\newcommand{\orcidauthorB}{0000-0002-8985-4728} % Add \orcidB{} behind the author's name

% Authors, for the paper (add full first names)
\Author{Eugen-Richard Ardelean $^{1}$\orcidA{}, and Raluca Portase $^{1,}$*\orcidB{}}

%\longauthorlist{yes}

% MDPI internal command: Authors, for metadata in PDF
\AuthorNames{Eugen-Richard Ardelean and Raluca Portase}

% Affiliations / Addresses (Add [1] after \address if there is only one affiliation.)
\address[1]{%
	$^{1}$ \quad Department of Computer Science, Technical University of Cluj-Napoca, Str. Memorandumului, Nr. 28, Cluj-Napoca, 400114, Cluj, Romania; richard.ardelean@cs.utcluj.ro\\; raluca.portase@cs.utcluj.ro
}

% Contact information of the corresponding author
\corres{Correspondence: raluca.portase@cs.utcluj.ro}

% Current address and/or shared authorship
%\firstnote{Current address: Affiliation.}  
% Current address should not be the same as any items in the Affiliation section.

%\secondnote{These authors contributed equally to this work.}
% The commands \thirdnote{} till \eighthnote{} are available for further notes.

%\simplesumm{} % Simple summary

%\conference{} % An extended version of a conference paper

% Abstract (Do not insert blank lines, i.e. \\) 
\abstract{Zero-Inflated Time Series (ZITS), characterised by excessive zeros relative to their standard distributions, pose significant challenges for synthetic data generation due to the need to simultaneously preserve the temporal dependencies, the excess zero structure, and the heavy-tailed non-zero distributions. Current generative models, including variational autoencoders (VAE) and generative adversarial networks (GAN), lack the mechanisms to handle zero-inflated data and fail to generate data with similar characteristics. This work proposes two architectures, ZITS-GAN and ZITS-VAE, designed for zero-inflated time series generation through a two-head decoder structure that separately models the zero pattern and the non-zero magnitude. The proposed architectures use dilated convolutional layers to capture long-sequence dependencies and stabilise training through a $\beta$-weighted ELBO for ZITS-VAE and a Wasserstein GAN with gradient penalty for ZITS-GAN. We evaluated the proposed models on two real-world time series datasets with zero ratios of 59\% and 74\% and demonstrated superior performance compared to state-of-the-art models, including TimeGAN, TimeVAE, ChronoGAN, TransFusion, and FIDE.}

% Keywords
\keyword{Zero-Inflated Data; Zero-Inflated Time Series; Time Series Generation; GAN; VAE; Synthetic Data}

% The fields PACS, MSC, and JEL may be left empty or commented out if not applicable
%\PACS{J0101}
%\MSC{}
%\JEL{}

%%%%%%%%%%%%%%%%%%%%%%%%%%%%%%%%%%%%%%%%%%
% Only for the journal Diversity
%\LSID{\url{http://}}

%%%%%%%%%%%%%%%%%%%%%%%%%%%%%%%%%%%%%%%%%%
% Only for the journal Applied Sciences
%\featuredapplication{Authors are encouraged to provide a concise description of the specific application or a potential application of the work. This section is not mandatory.}
%%%%%%%%%%%%%%%%%%%%%%%%%%%%%%%%%%%%%%%%%%

%%%%%%%%%%%%%%%%%%%%%%%%%%%%%%%%%%%%%%%%%%
% Only for the journal Data
%\dataset{DOI number or link to the deposited data set if the data set is published separately. If the data set shall be published as a supplement to this paper, this field will be filled by the journal editors. In this case, please submit the data set as a supplement.}
%\datasetlicense{License under which the data set is made available (CC0, CC-BY, CC-BY-SA, CC-BY-NC, etc.)}

%%%%%%%%%%%%%%%%%%%%%%%%%%%%%%%%%%%%%%%%%%
% Only for the journal BioTech, Fishes, Neuroimaging and Toxins
%\keycontribution{The breakthroughs or highlights of the manuscript. Authors can write one or two sentences to describe the most important part of the paper.}

%%%%%%%%%%%%%%%%%%%%%%%%%%%%%%%%%%%%%%%%%%
% Only for the journal Encyclopedia
%\encyclopediadef{For entry manuscripts only: please provide a brief overview of the entry title instead of an abstract.}

%%%%%%%%%%%%%%%%%%%%%%%%%%%%%%%%%%%%%%%%%%
% Different journals have different requirements. Please check the specific journal guidelines in the "Instructions for Authors" on the journal's official website.
%\addhighlights{yes}
%\renewcommand{\addhighlights}{%
%
%\noindent This is an optional section, the goal of which is to increase the discoverability and readability of the article via search engines and other scholars. Highlights should not be a copy of the abstract, but a simple overview allowing the reader to quickly and simply find out what the article is about and what can be cited from it. Each of these parts should be up to 2 bullet points.\vspace{3pt}\\
%\textbf{What are the main findings?}
% \begin{itemize}[labelsep=2.5mm,topsep=-3pt]
% \item First bullet.
% \item Second bullet.
% \end{itemize}\vspace{3pt}
%\textbf{What are the implications of the main findings?}
% \begin{itemize}[labelsep=2.5mm,topsep=-3pt]
% \item First bullet.
% \item Second bullet.
% \end{itemize}
%}

%%%%%%%%%%%%%%%%%%%%%%%%%%%%%%%%%%%%%%%%%%
\begin{document}

%%%%%%%%%%%%%%%%%%%%%%%%%%%%%%%%%%%%%%%%%%
\newcommand{\best}[1]{\textbf{#1}}
\newcommand{\scnd}[1]{\underline{#1}}

\section{Introduction}
A time series is a sequence of observations indexed by time \cite{ref1,refX} whose structure includes temporal dependencies, seasonality, and nonstationary behaviour. Foundational work \cite{ref1} on time series modelling and forecasting stressed the importance of capturing these dependencies and seasonality to understand temporal data. Modern data-driven applications, from sensors and energy to finance and biomedical data, rely on realistic time-series models for forecasting, prediction, inference and synthetic data generation for downstream tasks \cite{ref2,ref3}.

Generating synthetic time series has gained traction \cite{refNew} for model evaluation, benchmarking, and enabling data sharing when privacy or commercial limitations prevent distributing raw data. Many domains suffer from data shortages due to complex or noisy acquisition environments and even a scarcity of events \cite{tolas2024gemsyd, ref4}. Deep generative approaches such as TimeGAN \cite{ref2} have shown promising results in capturing complex temporal distributions \cite{ref5}, and more recent work integrates privacy-preserving mechanisms (e.g., differential privacy \cite{ref6}) to create synthetic series that preserve the characteristics of the real data while being suitable for sharing. An example is the generation of high-fidelity synthetic health data \cite{ref7,ref8}, which is particularly relevant given the high sensitivity of personal medical records.

Zero inflation \cite{ref9,ref10} is the phenomenon in which a time series contains an excessive number of zeros relative to the standard non-zero distributions; modelling this phenomenon has long been a statistical challenge \cite{ref9}. The canonical formulation specifies a two-component process: a structural-zero component and a conventional count component \cite{ref9}. The literature distinguishes between zero-inflated and hurdle formulations \cite{ref11,ref12,ref13}, with hurdle formulations having the advantage of not necessarily assuming that zeros and non-zeros come from the same data-generating process.

From the perspective of synthetic data generation, zero-inflated time series (ZITS, also referred to as sparse time series or intermittent time series) are especially hard because they introduce several difficulties: (i) the need to reproduce complex dependencies (autocorrelation, seasonality), (ii) the need to reproduce the excess zeros and the heavy tails in the non-zero distribution, and (iii) the need to respect structural zeros and conditional zero-generating mechanisms that may depend on various data-dependent factors. Deep neural networks excel at capturing soft temporal dependencies but can struggle to represent the rare-event structure found in ZITS \cite{ref14,ref15,ref16,ref17} without careful architectural choices, loss designs, or conditioning strategies. These limitations can be extended to deep generators such as  RNNs, GANs, VAEs, or hybrid architectures.

The two most popular approaches for handling ZITS are zero-inflated Poisson (ZIP) regression models \cite{ref9,ref19} and hurdle models \cite{ref19,ref20}. Both methods separate the data into zero and non-zero components \cite{ref19}, with ZIPs modelling excess zeros and hurdles modelling the entire distribution of zeros and the non-zeros of interest \cite{ref19}. In the literature, both methods have been applied to various domains, such as manufacturing defects \cite{ref9}, road crash frequency \cite{ref5}, insurance claims \cite{ref21}, and domestic violence \cite{ref22}.

The contributions of this work are:
\begin{itemize}
    \item Adapting the zero-inflated neural network framework \cite{ref19} originally developed for static insurance data to the temporal domain—moving beyond static mixture models or marginal distributions — and introducing an optional temporal consistency regularisation loss designed to assist in capturing realistic run-length statistics in zero and non-zero patterns where needed, and its effect is examined in the hyperparameter search (Appendix~\ref{app:hyperparams})).
    \item Providing a comparative evaluation of our proposed ZITS-VAE and ZITS-GAN models for zero-inflated time series generation against traditional and modern deep neural network methods \cite{ref2,ref23,ref4,ref29}, demonstrating higher performance in preserving both the zero-inflation structure and non-zero value distributions.
    \item Establishing a reusable framework that includes VAE and GAN variants and can be extended to other zero-inflated time series applications, where structural zeros represent a fundamental challenge for standard generative models.
\end{itemize}

\section{Materials and Methods}
\subsection{Related Work}
TimeVAE \cite{ref23} is a variational autoencoder (VAE) architecture designed for multivariate time series that prioritises both interpretability and training cost. TimeVAE uses the autoencoder to denoise and compress temporal signals, and its decoder submodule can inject interpretable trend and seasonality components into generated samples. TimeVAE optimises the evidence lower bound (reconstruction term plus a Kullback-Leibler penalty). It was shown to yield competitive performance based on similarity metrics while reducing training time compared to other models.

Two baseline VAE variants have been paired with TimeVAE \cite{ref23} in evaluations, specifically the dense (VAE-Dense \cite{ref23}) and convolutional (VAE-Conv \cite{ref23}) variants. VAE-Dense uses fully connected encoder/decoder blocks, while VAE-Conv replaces them with one-dimensional convolutional layers to capture temporal structure. Both VAE-Conv and VAE-Dense served as parameter-efficient baselines.

TimeGAN \cite{ref2} constructs a hybrid architecture that learns a low-dimensional embedding of the time series data and then adversarially matches that embedding space to generate realistic time series. It combines an embedder–recovery autoencoder to map between raw time series space and a learned latent space, a generator–discriminator adversarial pair to make generated embeddings indistinguishable from the embedded real data, and a supervisor that enforces temporal dependencies inside the latent space. Training optimises multiple losses jointly: an autoencoder reconstruction loss (mean squared error), a supervised stepwise prediction loss (mean squared error), an unsupervised adversarial loss (binary cross-entropy), and auxiliary moment-matching terms. Alongside the TimeGAN model \cite{ref2}, the authors also proposed two post-hoc evaluation (later widely adopted \cite{ref25}) metrics, the Discriminative Score (DS) and the Predictive Score (PS). 

BetaVAE \cite{ref26,ref27} modifies the standard VAE objective by multiplying the KL divergence term by a scalar factor $\beta$ > 1. This simple reweighting increases disentanglement in representations that match the real data, with higher $\beta$ values typically forcing simpler latent representations at the expense of reconstruction fidelity.

ChronoGAN \cite{ref4} frames generation as a joint autoencoder–GAN problem in which the autoencoder constructs a compact latent embedding and a GAN operates in that embedding while the discriminator provides feedback in the original feature space. ChronoGAN also uses a supervisory network that enforces stepwise conditional structure, and an “early generation” protocol that snapshots generator outputs during training to counter training instability using the Discriminative and Predictive scores \cite{ref2,ref25}, and a hybrid GRU–LSTM encoder/decoder backbone intended to generalise across short and long sequences. 

Frequency-Inflated Conditional Diffusion Model for Extreme-Aware Time Series Generation (FIDE) \cite{ref28} extends conditional diffusion modelling by explicitly manipulating frequency-domain representations during generation and incorporates the Generalised Extreme Value (GEV) distribution into its generative model. FIDE adapts score-matching with frequency-domain reshaping to better maintain rare events in the generated data; a property that is particularly relevant when synthetic data must reproduce heavy tails or rare spikes for applications in risk, climate, and certain biomedical signals.

TransFusion \cite{ref29} is a diffusion model that integrates transformer encoders into the backward denoising process to generate high-fidelity long sequences. TransFusion introduces two main new concepts: diffusion-based probabilistic generation to avoid adversarial training issues and transformer attention to capture long-range temporal dependencies, enabling generation of sequences significantly longer than those typically reported for GAN or RNN-based generators. To evaluate the quality of the generated data, we introduced two transformer-based evaluation metrics (Long-Sequence Discriminative Score and Long-Sequence Predictive Score) \cite{ref29} that replace the TimeGAN post-hoc classifier and predictor scores with transformer models better suited to long sequences.

TSGM \cite{ref30} is an extensible framework for time-series generative modelling, not a single model that implements TimeGAN and TimeVAE.


While deep generative models have advanced significantly for standard continuous time series, their underlying architectural and objective design renders them ill-equipped for zero-inflated time series without modification. Existing models can be critically evaluated across the following primary failures:


\begin{itemize}
    \item Reconstruction bias (TimeGAN, TimeVAE, VAE-Dense/Conv): These architectures rely on smooth continuous likelihoods or mean-squared error (MSE) reconstruction losses \cite{ref2, ref23}. In a zero-inflated distribution with high amounts of exact zero observations, minimising MSE penalises errors smoothly. Because the non-zero magnitude values occupy a wide dynamic range, the continuous loss gradient dominates, pulling predictions away from exact zero and resulting in a total failure to generate structural zeros (as demonstrated by a 0.0 zero ratio in our empirical evaluation). They lack a discrete-continuous hybrid likelihood or a dedicated gating mechanism.
    \item Mode collapse on sparsity (ChronoGAN): Adversarial frameworks struggle with highly skewed, bimodal or sparse distributions \cite{ref4}. When the training data is heavily dominated by zeros, the discriminator and generator dynamics can become unstable, frequently causing the generator to collapse into the opposite degenerate extreme—producing sequences consisting entirely of zeros.
    \item Lack of structural zero objectives (FIDE, TransFusion): Advanced models like FIDE \cite{ref28} (which targets extreme values via GEV and frequency-domain diffusion) and TransFusion \cite{ref29} (which leverages transformer-based diffusion for long sequences) are engineered to capture heavy tails and long-range dependencies. However, they assume smooth probability distributions. Consequently, they smooth out discrete zero states and fail to preserve the exact point-mass probability required for intermittent series.
    \item Absence of temporal intermittency modelling (ZIP and Hurdle Models): While classical statistical baselines \cite{ref9} successfully model marginal zero probabilities per timestep, they treat each timestep independently. By failing to model serial dependence across time, they cannot capture zero-clustering, leading to higher autocorrelation errors.
\end{itemize}



\subsection{Traditional Statistical Baselines for Zero-Inflated Time Series}
To situate our proposed models against classical statistical approaches, we include two well-established baselines for zero-inflated data: the Zero-Inflated Poisson (ZIP) model and the Hurdle (two-process) model. Both are fit independently at each timestep $t \in \{1, \ldots, T\}$, treating the $N$ observed values $\{x_{i,t}\}_{i=1}^{N}$ as i.i.d. draws from a marginal distribution. This per-timestep formulation means neither baseline captures temporal dependence across steps, which is an intentional design choice: they represent the strongest purely marginal statistical competitors, isolating the contribution of temporal modelling in our deep generative models. 


The Zero-Inflated Poisson model \cite{ref9} augments a standard Poisson distribution with a point mass at zero to account for an excess of zeros beyond what the Poisson rate alone would predict. For timestep $t$, the probability mass function is
\begin{equation}
  P(x_t = k) =
  \begin{cases}
    \pi_t + (1 - \pi_t)\,e^{-\lambda_t} & k = 0, \\
    (1 - \pi_t)\,\dfrac{\lambda_t^k\,e^{-\lambda_t}}{k!} & k = 1, 2, \ldots
  \end{cases}
\end{equation}
where $\pi_t \in (0,1)$ is the probability of a \emph{structural zero} (an observation that is zero by mechanism, independently of the count process) and $\lambda_t > 0$ is the Poisson rate of the \emph{sampling} component. The model therefore distinguishes between two types of zero: structural zeros generated with probability $\pi_t$, and \emph{sampling zeros} that arise from the Poisson component when the count happens to be zero.

We estimate the parameters $(\pi_t, \lambda_t)$ via the Expectation-Maximisation (EM) algorithm \cite{Dempster1977}. In the E-step, the posterior probability that the $i$-th observation is a structural zero is
\begin{equation}
  w_{i,t} =
  \begin{cases}
    \dfrac{\pi_t}{\pi_t + (1 - \pi_t)\,e^{-\lambda_t}} & x_{i,t} = 0, \\
    0 & x_{i,t} > 0.
  \end{cases}
\end{equation}

In the M-step, the parameters are updated as
\begin{align}
  \pi_t &\leftarrow \frac{1}{N}\sum_{i=1}^{N} w_{i,t}, \\
  \lambda_t &\leftarrow \frac{\sum_{i=1}^{N} x_{i,t}(1 - w_{i,t})}
                            {\sum_{i=1}^{N} (1 - w_{i,t})}.
\end{align}

The E- and M-steps are iterated until $|\pi_t^{(\ell+1)} - \pi_t^{(\ell)}|$ and $|\lambda_t^{(\ell+1)} - \lambda_t^{(\ell)}|$ both fall below $10^{-6}$ or 200 iterations are reached.

The Hurdle model \cite{Cragg1971, ref20} is a two-process model that separates the zero/non-zero decision from the distribution of positive values. Unlike the ZIP model, a hurdle model does not decompose zeros into structural and sampling types; instead, \emph{all} zeros are generated by a single binary gate, and the positive part is modelled independently. For timestep $t$:
\begin{align}
  P(x_t = 0) &= \pi_t, \\
  P(x_t \mid x_t > 0) &= f_{\text{pos}}(x_t;\,\theta_t),
\end{align}
where the gate probability $\pi_t \in (0,1)$ and the positive-part parameters $\theta_t$ are estimated separately. We model the positive component with a Log-Normal distribution,
\begin{equation}
  \ln x_t \mid x_t > 0 \;\sim\; \mathcal{N}(\mu_t,\, \sigma_t^2),
\end{equation}
which is well-suited for heavy-tailed, right-skewed positive data, and provides a reasonable approximation for count data when the support is large.

The gate probability is estimated by maximum likelihood as the sample proportion of zeros,
\begin{equation}
  \hat{\pi}_t = \frac{1}{N}\sum_{i=1}^{N} \mathbf{1}[x_{i,t} = 0].
\end{equation}

MLE estimates the Log-Normal parameters on the positive observations $\{x_{i,t} : x_{i,t} > 0\}$,
\begin{align}
  \hat{\mu}_t &= \frac{1}{N_t^{+}}\sum_{i:\,x_{i,t}>0} \ln x_{i,t}, \\
  \hat{\sigma}_t &= \sqrt{\frac{1}{N_t^{+}}\sum_{i:\,x_{i,t}>0}
      (\ln x_{i,t} - \hat{\mu}_t)^2},
\end{align}
where $N_t^{+}$ is the number of positive observations at timestep $t$.




\subsection{Proposed Approach}
Our proposed approach adapts the interpretable zero-inflated neural network framework \cite{ref19} from insurance data to time series generation. The output is decomposed into two separate components that are trained with distinct loss functions. We adopt an explicit two-head architecture in which the decoder/generator produces a gate component to determine whether each time step is zero and a magnitude component to determine the value of each time step.


\begin{equation}
    \text{gate\_logit} \xrightarrow{\sigma} p_{\text{gate}} \in (0,1) \quad [\text{P(non-zero) at each timestep}] 
\end{equation}

\begin{equation}
    \text{mag\_logit} \xrightarrow{\sigma} \text{mag} \in (0,1) \quad [\text{normalized magnitude} \mid \text{non-zero}]
\end{equation}

During training, we supervise these two heads independently: the gate is trained with binary cross-entropy (BCE) against a binary mask, while the magnitude is trained with mean squared error (MSE) restricted to non-zero entries. During generation, the final output is computed as:

\begin{equation}
    \text{binary}_t \sim \text{Bernoulli}(p_{\text{gate},t})
\end{equation}

\begin{equation}
    \hat{x}_t = \text{binary}_t \times \text{mag}_t
\end{equation}

As zeros in the data can cluster, an optional temporal consistency (TC) loss was introduced as the lag-1 autocorrelation, which penalizes mismatches in the autocorrelation structure of the zero/non-zero patterns while also encouraging the reproduction of statistics:

\begin{equation}
  \mathcal{L}_{\text{TC}} = \left( \mathbb{E}[p_{\text{gate},t} \cdot p_{\text{gate},t+1}] - \mathbb{E}[\mathbf{1}_{x_t>0} \cdot \mathbf{1}_{x_{t+1}>0}] \right)^2
\end{equation}

The TC term is optional: setting $\alpha_{tc}=0$ removes it, in which case temporal dependencies are captured only by the dilated convolutional layers. Whether the term is beneficial depends on the model and dataset, so $\alpha_{tc}$ is included in the hyperparameter search (Appendix~\ref{app:hyperparams}) rather than fixed a priori.


Both our proposed VAE and GAN use dilated convolutional layers \cite{ref31}. With exponentially increasing dilation rates (1, 2, 4), these layers expand the receptive field without reducing temporal resolution, and their convolutions can be computed in parallel. Residual connections are used to support gradient flow through the convolutional blocks. These architectural properties motivate our choice of dilated convolutions; however, we did not conduct a direct controlled comparison against an LSTM, and therefore do not claim empirical superiority over LSTMs. The M5 comparison between dilated and standard convolutions is reported in Appendix~\ref{app:ablation}.




Our VAE combines the standard evidence lower bound (ELBO) with the zero-inflated decomposition:

\begin{equation}
  \mathcal{L}_{\text{VAE}} = \beta \cdot \text{KL}(q(z \mid x) \parallel p(z)) + \alpha_r \cdot \mathcal{L}_{\text{recon}} + \alpha_g \cdot \mathcal{L}_{\text{gate}} + \alpha_{tc} \cdot \mathcal{L}_{\text{TC}}
\end{equation}

where:
\begin{itemize}
    \item $\mathcal{L}_{\text{recon}} = \text{MSE}(x_{\text{mask}}, \hat{x}_{\text{mask}})$ applied only to non-zero timesteps
    \item $\mathcal{L}_{\text{gate}} = \text{BCE}(\text{gate\_logit}, \mathbf{1}_{x>0})$ the sparsity predictor
    \item $\mathcal{L}_{\text{TC}}$ enforces temporal consistency in zero patterns
  \item $\beta$ controls the KL regularization strength (set to $0.3$).
\end{itemize}

The loss weights $\alpha_{r}, \alpha_{g}, \alpha_{tc}$ were selected by random search with Borda-count aggregation (Appendix~\ref{app:hyperparams}); the selected values are given in Table~\ref{tab:hyperparam}. 


Our GAN employs Wasserstein GAN with gradient penalty (WGAN-GP), which provides training stability for highly sparse distributions where standard GANs can collapse.

The discriminator uses spectral normalisation on all convolutional layers (eliminating batch normalisation, which is incompatible with WGAN-GP):
\begin{equation}
  \mathcal{L}_{D} = \mathbb{E}\big[D(x_{\text{fake}})\big] - \mathbb{E}\big[D(x_{\text{real}})\big] + \lambda_{\text{gp}}\cdot \mathrm{GP},
\end{equation}
where the gradient penalty enforces the 1-Lipschitz constraint:
\begin{equation}
  \mathrm{GP} = \mathbb{E}_{x}\big(\|\nabla_{x} D(x)\|_{2} - 1\big)^{2}, \qquad
  x = \epsilon x_{\text{real}} + (1-\epsilon) x_{\text{fake}}.
\end{equation}



The generator combines adversarial loss with a zero-inflation component:
\begin{equation}
  \mathcal{L}_{G} = -\mathbb{E}\big[D(x_{\text{fake}})\big]
  + \alpha_{r}\cdot \mathcal{L}_{\text{recon}}
  + \alpha_{g}\cdot \mathcal{L}_{\text{gate}}
  + \alpha_{fm}\cdot \mathcal{L}_{\text{FM}}
  + \alpha_{tc}\cdot \mathcal{L}_{\text{TC}}.
\end{equation}

The loss weights $\alpha_{r}, \alpha_{g}, \alpha_{fm}, \alpha_{tc}$ were selected by random search with Borda-count aggregation (Appendix~\ref{app:hyperparams}); the selected values are given in Table~\ref{tab:hyperparam}. 


Feature matching loss matches intermediate discriminator activations to stabilise training:
\begin{equation}
  \mathcal{L}_{\text{FM}} = \big\| \mathbb{E}\big[\phi(x_{\text{real}})\big] - \mathbb{E}\big[\phi(x_{\text{fake}})\big] \big\|_{2}^{2},
\end{equation}
where $\phi(\cdot)$ extracts features from the second convolutional layer of the discriminator.


\subsection{Metrics}
Maximum Mean Discrepancy distance (MMD) \cite{ref33,ref34,ref35} is a kernel-based two-sample statistic that measures the distance between distributions by comparing mean embeddings in a reproducing kernel Hilbert space. MMD is sensitive to discrepancies in distributions and has the advantage of being a nonparametric test that has been used to evaluate the quality of the generated data \cite{ref8}.

Predictive Score (PS) \cite{ref2} is a post-hoc evaluation based on a prediction task in which a predictor is trained on the generated data and evaluated on the real dataset using mean absolute error (MAE). Lower MAE indicates better information retention for prediction, meaning the generation is more robust.

The Discriminative Score (DS) \cite{ref2} is a post-hoc evaluation based on a classification task in which a standard classifier is trained to distinguish real from synthetic sequences. The discriminative score is the classification error, and as such, lower values are better. It operates on the assumption that if a classifier cannot reliably separate between the real and generated data, the synthetic samples are close to the true distribution. 

Long-sequence Discriminative Score (LDS) \cite{ref29} is a transformers-based variant of the DS that takes into consideration that classifier training can deteriorate when using long sequences. It has been argued that a transformer encoder better captures long-range dependencies and thus provides a more stringent test \cite{ref29}. 

For DS and LDS, because they are classification approaches, a value near 0.5 means the classifier cannot distinguish between real and generated data and thus corresponds to higher similarity. However, DS and LDS are defined as $|0.5 - \text{accuracy}|$ when classifying real versus generated data, meaning that 0 indicates ideal performance. 

Long-sequence Predictive Score (LPS)  \cite{ref29} is a variant of the PS that uses a transformer-based sequence predictor trained on synthetic data and evaluated on real data. The performance metric of LPS is MAE for one-step prediction.

For a comprehensive evaluation of the similarity between the real and generated data, simpler statistical similarity values have also been computed as the differences: between zero ratios, means, standard deviations,  various quantile values (for example, Q10, Q25,...), skewness and kurtosis, but also the Wasserstein distance, KL divergence and the Autocorrelation MAE.

\subsection{Data}
The first dataset was the M5 forecasting data \cite{ref38} provided on Kaggle, which contains daily unit sales collected by Walmart and distributed as several CSV files. The raw release includes a wide table of bottom-level series; each row corresponds to a single stock keeping unit (SKU) at a single store, along with a separate calendar mapping that links each recorded day index to its calendar date and auxiliary fields such as weekday, month, and special events. This publicly available data comprises 30,490 items over 1,941 days. The M5 data explicitly focuses on realistic retail demand and displays intermittency (lots of zero daily sales). All 30,490 samples of the M5 dataset have been used in the analyses.

Preprocessing converted the released tables into a two-dimensional matrix whose rows represent bottom-level series and whose columns represent consecutive calendar days covering one year. To construct this matrix, we aligned the daily sales history to calendar dates and selected a contiguous 365-day window ending on the most recent available date to capture a complete annual seasonal cycle while ensuring computational tractability. All entries are integer counts of units sold on the corresponding date; entries equal to zero indicate zero sales of that item on that date. The data statistics are as follows: zero ratio of 0.592, a mean of 1.311, a standard deviation of 3.634, a skewness of 11.503, and a kurtosis of 284.567.

To generate the data, we needed additional preprocessing. Because count data are also highly skewed and contain occasional extreme values, we used a robust logarithmic transformation of the form $log10(x+1)/s$. We determined the scale factor from the 99th percentile of the non-zero values rather than the absolute maximum, making the normalisation less sensitive to outliers. This transformation preserves zeros exactly while compressing large counts into a narrower dynamic range. The inverse transform recovers count values by exponentiation, clips to the maximum observed count, and rounds to the nearest integer.



The second dataset used for experimental validation consists of time series data capturing the daily operational duration of household IoT devices characterised by discrete operating cycles over one year. To ensure statistical relevance and data quality, we retained from the raw dataset - comprising individual running cycle durations in seconds per device - only those devices exhibiting a sufficient number of running cycles. The following filtering criteria were applied, resulting in a total of 8.1 million running cycles across more than 40,000 devices:
\begin{itemize}
    \item We retained only running cycles between 5 minutes and 12 hours, as values outside this range are implausible for the device types considered.
    \item Only devices with more than 20 running cycles per year were included, ensuring adequate statistical significance based on frequency of use (acknowledging that this filters out extremely infrequent edge cases).
\end{itemize}
After cleaning the dataset, we aggregated the running cycles for each device at a daily time resolution. In the resulting time series, each data point represents a device's total runtime (in seconds) on a given day of the year. The statistics are as follows: a zero ratio of 0.741, a mean of 2569, a standard deviation of 5633, a skewness of 2.851, and a kurtosis of 9.792. The household IoT devices dataset contains 49,208 samples, and all samples were used in the analyses.

To generate the data, we required additional preprocessing. To reduce the strong right skew typical of duration data, we applied a logarithmic normalisation of the form $log10(x/3600+1)$. We later reused the fitted normalisation for the inverse transform, as the newly generated data contained normalised values that we transformed back to seconds, clipping to the observed physical range.


\section{Results}
We evaluated ZITS generation on two real-world datasets. The M5 dataset \cite{ref38} presents retail sales data with a zero ratio of 0.592, while the household IoT devices dataset presents a more challenging task with a zero ratio of 0.741. In both cases, the classical statistical baselines, ZIP and Hurdle, provide an important point of comparison because they are explicitly designed for zero-inflated data.

Our proposed models used the following architectural choices and hyperparameters:
\begin{itemize}
    \item {Number of channels for convolutional layers:} 64.
    \item {Latent variables:} 64.
    \item {Activation function:} Leaky RELU (0.2).
    \item {Dropout:} 0.3.
    
    \item {Optimizer:} Adam (for the GAN variant: 5 discriminator steps per generator step). 
    \item {Batch size:} 64.
    \item {Gradient clipping:} Max norm 1.0 to prevent exploding gradients.
    \item {Early stopping:} Patience of 20 epochs on validation loss.
    \item {Other hyperparameters:} as given by Table \ref{tab:hyperparam} found through a random search across various parameter configurations (detailed in the Appendix, see section Appendix ~\ref{app:hyperparams} and Tables \ref{tab:hyperparams_iot_gan}, \ref{tab:hyperparams_iot_vae}, \ref{tab:hyperparams_m5_gan} and \ref{tab:hyperparams_m5_vae}). The best values were chosen through a Borda ranking from 50 configurations that were trained for 80 epochs.
\end{itemize}

\begin{table}[htbp]
\centering
\caption{Best parameters for the models obtained through a Random Search across various parameter configurations (detailed in the Appendix; see section~\ref{app:hyperparams})}
\label{tab:hyperparam}
\begin{tabular}{l cccc}
\toprule
& \multicolumn{2}{c}{\textbf{IoT Data}} & \multicolumn{2}{c}{\textbf{M5 Data}} \\
\cmidrule(lr){2-3} \cmidrule(lr){4-5}
\textbf{Parameter} & \textbf{ZITS-GAN} & \textbf{ZITS-VAE} & \textbf{ZITS-GAN} & \textbf{ZITS-VAE} \\
\midrule
$\mathit{latent\_dim}$ & 128     & 128     & 32      & 128     \\
$\mathit{lr}$         & 0.0002  & 0.0005  & 0.0001  & 0.0005  \\
$\beta_0$    & 0.5     & 0.1     & 0.000   & 1       \\
$\alpha_{g}$ & 1       & 10      & 1       & 5       \\
$\alpha_{r}$ & 10      & 10      & 10      & 5       \\
$\alpha_{tc}$    & 0.000   & 0.5     & 0.000   & 0       \\
$\alpha_{fm}$    & 0.5     & ---     & 2       & ---     \\
$\mathit{train time (s)}$   & 322.8   & 255.7   & 242.7   & 370.6   \\
\bottomrule
\end{tabular}
\end{table}


Our proposed ZITS-VAE and ZITS-GAN models were compared against state-of-the-art models including TimeGAN, TimeVAE, ChronoGAN, TransFusion, FIDE, and several VAE variants, as well as traditional statistical baselines (ZIP and Hurdle) to assess their ability to preserve zero-inflation structure, statistical properties, and temporal dependencies.  

Table~\ref{tab:m5_results} shows the results on the M5 dataset. The most immediate and consequential finding concerns the zero ratio: the vast majority of state-of-the-art deep generative models — FIDE, TimeVAE, TimeGAN, TransFusion, and the VAE variants — fail to produce \emph{any} zeros, generating only positive values. This is a failure case: a synthetic dataset with no zeros cannot faithfully represent a retail distribution where nearly 60\% of daily sales are zero. ChronoGAN exhibits the opposite degenerate behaviour, generating \emph{only} zeros, which is equally unusable. TSGM-TimeVAE partially recovers sparsity (zero ratio difference of 0.144; synthetic zero ratio of 0.736) but overcorrects beyond the true ratio.

By contrast, ZITS-GAN and ZITS-VAE achieve zero-ratio differences of only 0.012 and 0.022, respectively, because the explicit two-head decoder separates the binary gating decision from magnitude estimation, preventing the magnitude head from dominating and suppressing all zeros. This collapse occurs in single-head architectures when the non-zero reconstruction loss overwhelms any implicit sparsity signal. The traditional statistical baselines (ZIP and Hurdle) also reproduce the zero ratio exactly (difference of 0.000) because they directly model the marginal zero probability at each timestep; however, they do so without temporal modelling, which matters when assessing autocorrelation fidelity. ZITS-VAE achieves the best autocorrelation MAE (0.046), followed by ZITS-GAN (0.063), while most other models — including TimeGAN (0.296), TimeVAE (0.405), and TransFusion (0.348) — incur substantially larger errors. The statistical baselines (ZIP: 0.081, Hurdle: 0.082) perform better than most deep models but are surpassed by both ZITS variants. This occurs because ZIP and Hurdle are fit independently at each timestep and thus produce serially uncorrelated samples.

For distributional shape, ZITS-VAE achieves the best mean difference (0.300) and the best standard deviation difference (1.296) among all deep generative models. This advantage reflects that the magnitude head is trained exclusively on non-zero entries, avoiding the zero-bias that degrades mean and variance estimates in standard architectures where zero and non-zero values compete within the same reconstruction loss. Higher-order moments remain challenging for all models: skewness and kurtosis differ widely across the board, attributable to the M5 distribution's extreme heavy-tailedness (true kurtosis of 284.567). Nevertheless, ZITS-GAN achieves the lowest kurtosis difference (277.47) among deep models, reflecting that the Wasserstein adversarial objective, with its emphasis on distribution-level matching rather than pointwise reconstruction, better approximates extreme quantile behaviour than reconstruction-based objectives.

Quantile preservation reinforces this picture. Both ZITS models achieve perfect matching at the 10th, 25th, and 50th percentiles, which is consistent with their zero ratio: because a zero ratio of approximately 59\% implies that the median and lower quantiles lie at or near zero, any model that correctly generates sparsity will automatically match these quantiles without additional effort. At the 75th and 90th percentiles — where the distribution transitions into its non-zero body — ZITS-GAN achieves differences of 0.000 and 1.000, outperforming the statistical baselines, which accrue quantile errors in the non-zero range precisely because they model positive values independently of temporal context. At the 99th percentile (extreme tail), ZITS-GAN (difference of 11.000) and ZITS-VAE (5.000) remain competitive with the Hurdle model (4.641), indicating that the magnitude head captures the tail reasonably well even without an explicit heavy-tail distribution assumption such as the Log-Normal used by Hurdle.

For distributional distance measures, MMD is 0.002 for both ZITS models, matching ZIP, Hurdle, and the best competing deep model (FIDE). The Wasserstein distance (ZITS-GAN: 1.503, ZITS-VAE: 1.545) is competitive with the Hurdle model (0.728) and better than all other deep models; the residual gap relative to Hurdle is attributable to the fact that Hurdle's Log-Normal positive component directly models the distributional shape of non-zero values without having to learn it from data, giving it an advantage in distributional distance at the cost of ignoring all temporal structure. The best KL divergence scores are FIDE (0.010) and the Hurdle model (0.022), substantially outperforming all other deep models; however, FIDE generates no zeros, making its KL divergence misleadingly low because it concentrates mass away from zero.




\begin{table}[htbp]
\centering
\caption{Results on M5 data with a 0 ratio of $\sim$0.59. Rows represent models and columns represent performance metrics. Bold values indicate the best for each metric; underlined values indicate the second-best.}
\label{tab:m5_results}
\setlength{\tabcolsep}{2.2pt} 
\renewcommand{\arraystretch}{1.0} 
\begin{tabular}{l cccccccc}
\toprule
& \multicolumn{5}{c}{\textbf{Statistical Differences ($\Delta$)}} & & & \\
\cmidrule(lr){2-6}
\textbf{Model} & \textbf{0-rat} & \textbf{mean} & \textbf{std} & \textbf{skew} & \textbf{kurt} & & & \\
\midrule
\textbf{ZIP}          & \textbf{0.000} & \textbf{0.000} & 1.726 & 10.135 & 283.43 & & & \\
\textbf{Hurdle}       & \textbf{0.000} & 0.133 & \underline{1.331} & \textbf{6.944} & \underline{235.70} & & & \\
\textbf{ChronoGAN}    & 0.408 & 1.311 & 3.634 & nan & nan & & & \\
\textbf{FIDE}         & 0.592 & 199.76 & 709.88 & 14.045 & 1085.6 & & & \\
\textbf{TimeGAN}      & 0.592 & 0.494 & 3.295 & 11.47 & 283.3 & & & \\
\textbf{TimeVAE}      & 0.592 & 0.493 & 3.224 & \underline{7.486} & 74.48 & & & \\
\textbf{TransFusion}  & 0.592 & 0.671 & 3.434 & 10.90 & 285.4 & & & \\
\textbf{TSGM-TimeGAN} & 0.592 & 1.295 & 3.616 & 3.841 & \textbf{33.86} & & & \\
\textbf{TSGM-TimeVAE} & 0.144 & 1.311 & 3.634 & 1875.0 & 4.2e06 & & & \\
\textbf{VAE-Conv}     & 0.592 & 0.359 & 3.314 & 11.45 & 283.4 & & & \\
\textbf{VAE-Dense}    & 0.592 & 0.359 & 3.314 & 11.45 & 283.4 & & & \\
\textbf{ZITS-GAN}     & \underline{0.012} & 0.606 & 2.601 & 9.542 & 277.47 & & & \\
\textbf{ZITS-VAE}     & 0.022 & \underline{0.300} & \textbf{1.296} & 9.542 & 280.24 & & & \\
\midrule
& \textbf{Qmean} & \multicolumn{7}{c}{\textbf{Quantile Differences ($\Delta$Q)}} \\
\cmidrule(lr){3-9}
\textbf{Model} & \textbf{error} & \textbf{Q10} & \textbf{Q25} & \textbf{Q50} & \textbf{Q75} & \textbf{Q90} & \textbf{Q95} & \textbf{Q99} \\
\midrule
\textbf{ZIP}          & 1.714 & \textbf{0.000} & \textbf{0.000} & \textbf{0.000} & 2.000 & 1.000 & \underline{1.000} & 8.000 \\
\textbf{Hurdle}       & \textbf{0.938} & \textbf{0.000} & \textbf{0.000} & \textbf{0.000} & 0.634 & 0.618 & 0.670 & \underline{4.641} \\
\textbf{ChronoGAN}    & 3.571 & \textbf{0.000} & \textbf{0.000} & \textbf{0.000} & 1.000 & 3.000 & 6.000 & 15.00 \\
\textbf{FIDE}         & 508.29 & 7.327 & 24.73 & 72.59 & 151.05 & 413.32 & 730.37 & 2158.6 \\
\textbf{TimeGAN}      & 3.215 & 1.403 & 1.558 & 1.758 & 1.058 & 0.697 & 3.601 & 12.42 \\
\textbf{TimeVAE}      & 3.455 & 0.795 & 0.795 & 0.795 & 0.205 & 2.203 & 5.201 & 14.18 \\
\textbf{TransFusion}  & 3.242 & 0.417 & 0.475 & 0.585 & 0.218 & 2.000 & 5.000 & 14.00 \\
\textbf{TSGM-TimeGAN} & 3.566 & 0.015 & 0.015 & 0.015 & 0.985 & 2.985 & 5.985 & 14.96 \\
\textbf{TSGM-TimeVAE} & 3.571 & \textbf{0.000} & \textbf{0.000} & \textbf{0.000} & 1.000 & 3.000 & 6.000 & 15.00 \\
\textbf{VAE-Conv}     & 3.215 & 1.286 & 1.436 & 1.629 & 0.909 & 0.859 & 3.772 & 12.61 \\
\textbf{VAE-Dense}    & 3.215 & 1.286 & 1.436 & 1.629 & 0.909 & 0.859 & 3.772 & 12.61 \\
\textbf{ZITS-GAN}     & 2.143 & \textbf{0.000} & \textbf{0.000} & \textbf{0.000} & \textbf{0.000} & \underline{1.000} & 3.000 & 11.000 \\
\textbf{ZITS-VAE}     & \underline{1.286} & \textbf{0.000} & \textbf{0.000} & \textbf{0.000} & 2.000 & 2.000 & \textbf{0.000} & \textbf{5.000} \\
\midrule
& \multicolumn{2}{c}{\textbf{Distances}} & \textbf{KL} & \textbf{MMD} & \textbf{PS} & \textbf{DS} & \textbf{LPS} & \textbf{LDS} \\
\cmidrule(lr){2-3}
\textbf{Model} & \textbf{Auto} & \textbf{Wass} & \textbf{div} & \textbf{score} & \textbf{score} & \textbf{score} & \textbf{score} & \textbf{score} \\
\midrule
\textbf{ZIP}          & \underline{0.081} & 1.468 & 0.519 & \textbf{0.002} & 1.422 & \textbf{0.000} & \underline{1.241} & 0.292 \\
\textbf{Hurdle}       & 0.082 & \textbf{0.728} & \underline{0.022} & \textbf{0.002} & \textbf{1.336} & \textbf{0.000} & 1.243 & 0.372 \\
\textbf{ChronoGAN}    & nan & nan & nan & 1.000 & 1.368 & \textbf{0.000} & 1.388 & 0.343 \\
\textbf{FIDE}         & 0.324 & 207.252 & \textbf{0.010} & \textbf{0.002} & 1.359 & 0.324 & \textbf{1.198} & 0.483 \\
\textbf{TimeGAN}      & 0.296 & 1.855 & 3.488 & 1.001 & 1.392 & 0.010 & 1.469 & 0.448 \\
\textbf{TimeVAE}      & 0.405 & 2.396 & 0.974 & 1.001 & 1.438 & 0.159 & 1.482 & 0.331 \\
\textbf{TransFusion}  & 0.348 & 2.574 & 2.701 & 0.019 & 1.618 & \textbf{0.000} & 1.609 & \textbf{0.100} \\
\textbf{TSGM-TimeGAN} & 0.896 & 3.198 & 3.476 & 1.001 & 1.374 & 0.009 & 1.343 & 0.368 \\
\textbf{TSGM-TimeVAE} & 0.138 & 3.214 & 3.343 & 0.997 & \underline{1.356} & \textbf{0.000} & 1.344 & 0.389 \\
\textbf{VAE-Conv}     & 0.289 & 1.885 & 2.783 & 1.001 & 1.385 & \textbf{0.000} & 1.459 & 0.385 \\
\textbf{VAE-Dense}    & 0.289 & 1.885 & 2.783 & 1.001 & 1.440 & \textbf{0.000} & 1.467 & 0.494 \\
\textbf{ZITS-GAN}     & 0.063 & \underline{1.503} & 0.908 & \textbf{0.002} & 1.372 & \textbf{0.000} & 1.348 & \underline{0.297} \\
\textbf{ZITS-VAE}     & \textbf{0.046} & 1.545 & 0.222 & \textbf{0.002} & 1.370 & \textbf{0.000} & 1.341 & 0.396 \\
\bottomrule
\end{tabular}
\end{table}


Table~\ref {tab:iot_results} shows the results on the household IoT devices dataset. With a zero ratio of 74.1\%, this dataset presents a substantially harder generation problem: a model must leave three out of every four timesteps at exactly zero while still reproducing the heavy-tailed distribution of operating durations on active days. The failure modes observed on M5 are even more pronounced here. ChronoGAN, TimeVAE, TimeGAN, and both standalone VAE variants produce no zeros, indicating that higher sparsity makes it even harder for single-head architectures to resist the pull of reconstruction loss toward generating non-zero values everywhere. Conversely, TransFusion and TSGM-TimeVAE collapse to the opposite extreme, generating \emph{only} zeros. This bimodal failure pattern — either all zeros or no zeros — illustrates the instability of models. Without an explicit gating mechanism, generative models operating on sparse distributions tend to converge to one of the two degenerate fixed points, depending on the balance between the reconstruction signal (which pushes toward non-zero mean values) and the global sparsity of the training data (which can force gradients toward the zero solution when most target values are zero).

ZITS-GAN and ZITS-VAE are the only deep generative models to avoid both failure modes, with zero ratio differences of 0.024 and 0.015, respectively. The statistical baselines also recover the zero ratio exactly, confirming that this is tractable with explicit zero modelling; however, the Hurdle model collapses on several distributional metrics (e.g., Wasserstein distance of 0.0 due to a degenerate fit on the positive component given the extreme sparsity) and cannot be considered a reliable baseline at this zero ratio.

For distributional shape, ZITS-VAE achieves mean and standard deviation differences of 617.524 and 795.339, substantially better than all other deep models that manage to generate non-zero values (do note that TimeGAN has a 12.87 mean difference but a 3112 std difference; TimeGAN's apparently low mean difference is misleading, as it generates no zeros and thus misrepresents the true distribution). The skewness and kurtosis differences for ZITS-VAE (0.313 and 1.926) are the best among all models, including the statistical baselines. This suggests that on distributions with moderate rather than extreme excess kurtosis (true kurtosis of 9.792 on IoT, versus 284.567 on M5), the ZITS magnitude head learns the non-zero distribution shape more faithfully than on M5, where extreme outliers make distributional matching harder for any parameterised model.

Quantile preservation follows the same pattern as M5: both ZITS models achieve perfect matching at the 10th, 25th, and 50th percentiles, again a direct consequence of a close zero ratio between the original and generated data. At the 75th and higher percentiles, ZITS-VAE (differences of 2830, 2162, 2221, 2246) and ZITS-GAN (differences of 3569, 2822, 1546, 3425) both substantially outperform all other deep models, though they do not match the Hurdle model at the 75th percentile, where Hurdle achieves a difference of only 4.991. This advantage at specific quantiles is again attributable to the Hurdle model's parametric assumption of a Log-Normal positive component; for operating duration data, the Log-Normal distribution reasonably approximates the marginal positive distribution, giving Hurdle an edge at individual quantiles even though it lacks temporal modelling.

Autocorrelation MAE demonstrates the strongest differentiation between model classes. ZITS-VAE achieves an autocorrelation MAE of 0.080 and ZITS-GAN of 0.106, both substantially lower than ChronoGAN (0.885), TimeGAN (0.254), TimeVAE (0.643), and TSGM-TimeGAN (0.877). Notably, the Hurdle model achieves 0.0, an artefact of the degenerate fit noted above: when the Hurdle model generates data with extremely low variance (collapsing to near-zero values), the resulting autocorrelation structure trivially matches the real data's binary pattern. 

For distributional distance metrics, ZITS-VAE achieves the lowest KL divergence (0.115) among all models, and the lowest Wasserstein distance (1695.392) among deep generative models. The MMD of 0.002 matches the best models. The Long-Sequence Discriminative Score for ZITS-GAN is 0.250, and for ZITS-VAE it is 0.060; the latter is the best among all deep models, indicating that the VAE variant generates sequences whose temporal structure is most difficult for a transformer-based classifier to distinguish from real data. The Long-Sequence Predictive Score for both ZITS models (approximately 2377) substantially outperforms all other models (which obtain values near 2847), confirming that sequences generated by the proposed models preserve the temporal dependencies relevant for one-step-ahead prediction tasks.


% Add \usepackage{booktabs} to your document preamble if not already present



\begin{table}[htbp]
\centering
\caption{Results on household IoT devices data with a 0 ratio of $\sim$0.74. Bold values indicate the best for each metric; underlined values indicate the second-best.}
\label{tab:iot_results}
\setlength{\tabcolsep}{2.2pt}
\renewcommand{\arraystretch}{1.0}
\begin{tabular}{l cccccccc}
\toprule
& \multicolumn{5}{c}{\textbf{Statistical Differences ($\Delta$)}} & & & \\
\cmidrule(lr){2-6}
\textbf{Model} & \textbf{0-rat} & \textbf{mean} & \textbf{std} & \textbf{skew} & \textbf{kurt} & & & \\
\midrule
\textbf{ZIP}          & 0.000 & \textbf{0.875} & 1270.590 & 1.725 & 10.482 & & & \\
\textbf{Hurdle}       & 0.000 & \underline{126.161} & 1072.073 & 2.125 & 39.903 & & & \\
\textbf{ChronoGAN}    & 0.741 & 33347 & 5480 & 5.055 & 58.87 & & & \\
\textbf{FIDE}         & 0.738 & 5.4e15 & 2.9e17 & 46.332 & 6690 & & & \\
\textbf{TimeGAN}      & 0.741 & 12.87 & 3112 & \underline{1.348} & 6.842 & & & \\
\textbf{TimeVAE}      & 0.741 & 2066 & 5365 & 16.09 & 348.0 & & & \\
\textbf{TransFusion}  & 0.259 & 2569 & 5633 & 1789 & 3.3e06 & & & \\
\textbf{TSGM-TimeGAN} & 0.741 & 2569 & 5633 & 11.268 & 203.81 & & & \\
\textbf{TSGM-TimeVAE} & 0.259 & 2569 & 5633 & nan & nan & & & \\
\textbf{VAE-Conv}     & 0.741 & 4.658 & 3288 & 1.4 & 7.194 & & & \\
\textbf{VAE-Dense}    & 0.741 & 19.74 & 3185 & 1.087 & \underline{5.104} & & & \\
\textbf{ZITS-GAN}     & 0.024 & 711.912 & \underline{910.615} & \underline{0.540} & 13.079 & & & \\
\textbf{ZITS-VAE}     & \underline{0.015} & 617.524 & \textbf{795.339} & \textbf{0.313} & \textbf{1.926} & & & \\
\midrule
& \textbf{Qmean} & \multicolumn{7}{c}{\textbf{Quantile Differences ($\Delta$Q)}} \\
\cmidrule(lr){3-9}
\textbf{Model} & \textbf{error} & \textbf{Q10} & \textbf{Q25} & \textbf{Q50} & \textbf{Q75} & \textbf{Q90} & \textbf{Q95} & \textbf{Q99} \\
\midrule
\textbf{ZIP}          & 3717.143 & \textbf{0.000} & \textbf{0.000} & \textbf{0.000} & 7558.000 & \textbf{116.000} & 3958.000 & 14388.000 \\
\textbf{Hurdle}       & \textbf{903.867} & \textbf{0.000} & \textbf{0.000} & \textbf{0.000} & \textbf{495.366} & \underline{470.329} & \textbf{300.304} & 5061.068 \\
\textbf{ChronoGAN}    & 28628 & 35890 & 35890 & 35890 & 34570 & 25840 & 21137 & 11183 \\
\textbf{FIDE}         & 2.2e13 & 17802 & 7643 & 6149 & 28394 & 4.0e07 & 5.9e09 & 1.6e14 \\
\textbf{TimeGAN}      & 4431 & 47.58 & 620.9 & 1814 & 2494 & 4031 & 7251 & 14758 \\
\textbf{TimeVAE}      & 7330 & 488.2 & 488.2 & 488.2 & \underline{831.7} & 9563 & 14302 & 25149 \\
\textbf{TransFusion}  & 7400 & \textbf{0.000} & \textbf{0.000} & \textbf{0.000} & 1320 & 10051 & 14791 & 25638 \\
\textbf{TSGM-TimeGAN} & 7400 & 0.017 & 0.017 & 0.017 & 1320 & 10051 & 14791 & 25638 \\
\textbf{TSGM-TimeVAE} & 7400 & \textbf{0.000} & \textbf{0.000} & \textbf{0.000} & 1320 & 10051 & 14791 & 25638 \\
\textbf{VAE-Conv}     & 4619 & 155.9 & 830.7 & 1937 & 2322 & 4197 & 7511 & 15380 \\
\textbf{VAE-Dense}    & 4485 & 193.7 & 736.7 & 1892 & 2296 & 4289 & 7479 & 14507 \\
\textbf{ZITS-GAN}     & 1623.040 & \textbf{0.000} & \textbf{0.000} & \textbf{0.000} & 3568.698 & 2822.448 & \underline{1545.588} & \underline{3424.546} \\
\textbf{ZITS-VAE}     & \underline{1351.279} & \textbf{0.000} & \textbf{0.000} & \textbf{0.000} & 2830.178 & 2161.948 & 2221.229 & \textbf{2245.594} \\
\midrule
& \multicolumn{2}{c}{\textbf{Distances}} & \textbf{KL} & \textbf{MMD} & \textbf{PS} & \textbf{DS} & \textbf{LPS} & \textbf{LDS} \\
\cmidrule(lr){2-3}
\textbf{Model} & \textbf{Auto} & \textbf{Wass} & \textbf{div} & \textbf{score} & \textbf{score} & \textbf{score} & \textbf{score} & \textbf{score} \\
\midrule
\textbf{ZIP}          & 0.114 & 4697.777 & 18.669 & \textbf{0.002} & 2699.790 & \textbf{0.000} & 2699.660 & 0.331 \\
\textbf{Hurdle}       & 0.116 & 1077.212 & \textbf{0.018} & \textbf{0.002} & 2699.640 & \textbf{0.000} & 2699.811 & 0.318 \\
\textbf{ChronoGAN}    & 0.885 & 26047 & 36.67 & 0.902 & 2677 & \textbf{0.000} & 2847 & 0.498 \\
\textbf{FIDE}         & \underline{0.089} & 1.2e16 & \textbf{0.001} & \textbf{0.002} & 2788 & \textbf{0.000} & 2847 & 0.497 \\
\textbf{TimeGAN}      & 0.254 & 7101 & 2.387 & \textbf{0.002} & 2578 & \textbf{0.000} & 2847 & 0.416 \\
\textbf{TimeVAE}      & 0.643 & 9418 & 31.029 & 0.898 & 2570 & 0.5 & 2847 & 0.469 \\
\textbf{TransFusion}  & 0.117 & 9919 & 20.166 & 1.001 & 2847 & \textbf{0.000} & 2847 & 0.5 \\
\textbf{TSGM-TimeGAN} & 0.877 & 9920 & 34.74 & 1.001 & 2677 & 0.276 & 2847 & 0.5 \\
\textbf{TSGM-TimeVAE} & nan & nan & nan & 1.001 & 2677 & \textbf{0.000} & 2847 & 0.5 \\
\textbf{VAE-Conv}     & 0.133 & 7322 & 3.472 & \textbf{0.002} & 2578 & \textbf{0.000} & 2847 & 0.468 \\
\textbf{VAE-Dense}    & 0.134 & 7168 & 2.086 & 0.021 & 2578 & 0.267 & 2847 & 0.4 \\
\textbf{ZITS-GAN}     & \textbf{0.062} & \underline{2210.168} & 0.155 & \textbf{0.002} & \textbf{2376.839} & \textbf{0.000} & \textbf{2376.940} & 0.250 \\
\textbf{ZITS-VAE}     & 0.080 & \textbf{1695.392} & \underline{0.115} & \textbf{0.002} & \underline{2376.919} & \textbf{0.000} & 2377.067 & \textbf{0.060} \\
\bottomrule
\end{tabular}
\end{table}
Our previous statistical analysis is complemented by our visual inspection of the overlap between the real and synthetically generated data through the use of the UMAP \cite{ref37} algorithm. Figure \ref{fig:m5_umap} shows 2-dimensional scatter plots of the entire M5 real dataset and the 30000 synthetically generated samples by each model processed through UMAP, while Figure \ref{fig:iot_umap} shows the 2-dimensional scatter plots of the entire household IoT devices dataset and 30000 synthetically generated samples. 

The UMAP embeddings of ZITS-GAN and ZITS-VAE substantially overlap with the real data, confirming that they explore and learn from the correct region of the data manifold. While UMAP projections are sensitive to hyperparameter choices and primarily reflect local neighborhood structures rather than absolute manifold coverage, the substantial visual overlap exhibited by ZITS-GAN and ZITS-VAE provides qualitative corroboration of their distributional alignment with the real data.



\begin{figure}[htbp]
  \centering
  \subfloat[UMAP scatter plots on the M5 sales dataset.\label{fig:m5_umap}]{%
    \includegraphics[width=0.7\textwidth]{fig/m5_umap.png}%
  }\\[2ex] % Adds a vertical space between the two subfigures
  \subfloat[UMAP scatter plots on the household IoT devices dataset.\label{fig:iot_umap}]{%
    \includegraphics[width=0.7\textwidth]{fig/iot_umap.png}%
  }%
  \caption{UMAP visualizations of original (red) and synthetically generated (blue) examples for (a) M5 and (b) household IoT datasets.}
  \label{fig:umap-stacked}
\end{figure}

The visual inspections of individual generated sequences (Figures~\ref{fig:m5-zits-gan} and~\ref{fig:m5-zits-vae} for the M5 dataset and Figures ~\ref{fig:iot-zits-gan} and~\ref{fig:iot-zits-vae} for the IoT dataset) further confirm that both models generate realistic intermittent patterns, with zero runs of plausible length interspersed with non-zero bursts of appropriate magnitude. We omit the sample plots for the other deep generative models, as they exhibit either constant zero or constant non-zero outputs, making visual comparison uninformative.

\begin{figure}[htbp]
  \centering
  \subfloat[ZITS-GAN.\label{fig:m5-zits-gan}]{%
    \includegraphics[width=0.48\textwidth]{fig/zits-gan-m5.png}%
  }%
  \hfill
  \subfloat[ZITS-VAE.\label{fig:m5-zits-vae}]{%
    \includegraphics[width=0.48\textwidth]{fig/zits-vae-m5.png}%
  }%
  \caption{Real (blue) and synthetically generated (orange) examples of the (a) ZITS-GAN and (b) ZITS-VAE proposed models for the M5 dataset.}
  \label{fig:zits-side-by-side}
\end{figure}

The NaN values that appear in certain cells arise because some models generate entirely constant sequences (all-zero or all-positive), which renders metrics such as skewness, kurtosis, Wasserstein distance, KL divergence, and autocorrelation undefined or infinite.


\section{Conclusion}
Our work addressed the challenge of generating synthetic zero-inflated time series (ZITS), a problem that has received limited attention despite its practical importance in domains such as retail demand, IoT device monitoring, insurance claims, and healthcare. Two architectures are proposed, ZITS-GAN and ZITS-VAE, that adapt the zero-inflated neural network framework to the temporal generation setting through two key design decisions and one optional component: (i) an explicit two-head decoder that separates zero-generation (gate head) from magnitude estimation (magnitude head); (ii) dilated convolutional layers with residual connections to capture long-range temporal dependencies without the training instabilities of recurrent architectures; and (iii) an optional temporal consistency loss that penalises mismatches in the lag-1 structure of the binary zero/non-zero pattern.

The fundamental motivation for the two-head design is that zero-inflated distributions are a mixture of two data-generating processes — a structural zero mechanism and a continuous positive distribution — and conflating these within a single output head inevitably causes one process to dominate during training. On both evaluation datasets (M5 retail, 59.2\% zeros; household IoT, 74.1\% zeros), every competing deep generative model under standard default tuning regimens without explicit zero handling collapsed to a degenerate solution due to class imbalance: either generating no zeros at all (because the reconstruction loss on non-zero values overwhelms any implicit sparsity signal) or generating only zeros (because the dominant zero mass in the training data drives the model to the trivial fixed point). 

All baseline models were executed using their official implementations and recommended hyperparameter settings adapted to our sequence lengths. Single-head models fail (producing either all zeros or no zeros) because standard reconstruction losses (MSE/BCE) on zero-inflated data are dominated by the majority class (zeros in sparse series or positive values in dense ones), causing gradient collapse toward a fixed point. 

ZITS-GAN and ZITS-VAE's ability to preserve zero-inflation structure improves on existing approaches that lack explicit mechanisms for handling ZITS. ZITS-GAN and ZITS-VAE were the only models to avoid both failure modes, achieving zero-ratio differences of at most 0.024 across both datasets, compared to deviations of 0.259 to 0.741 for most baselines. The statistical baselines (ZIP and Hurdle) match the zero ratio exactly — as they directly estimate the marginal zero probability — but produce serially uncorrelated samples, since they are fit independently at each timestep. ZITS-VAE and ZITS-GAN outperformed the statistical baselines on autocorrelation MAE on the M5 dataset (0.046 and 0.063 versus 0.081 and 0.082 for ZIP and Hurdle, respectively), indicating that the proposed architecture better reproduces the serial structure of zero clusters than per-timestep statistical models. We selected the TC loss with a non-zero weight in only one of the four final models, yet non-zero TC values always ranked in the top; the hyperparameter search suggests it may reduce autocorrelation error in some configurations, but we leave a controlled ablation for future work.

Distributional shape matching remains the most challenging aspect of ZITS generation, particularly for heavy-tailed distributions such as the M5 retail series (true kurtosis of 284.567). While ZITS models substantially outperform competing deep generators on skewness and kurtosis, they do not fully close the gap with parametric methods such as the Hurdle model with a Log-Normal positive component, which has a structural advantage when the assumed parametric form is a reasonable approximation of the true marginal distribution. This suggests a promising direction for future work: incorporating explicit heavy-tail distribution assumptions (e.g., Log-Normal, Pareto, or Generalised Extreme Value) into the magnitude head, similar to the approach taken by FIDE for non-zero values.

The proposed approaches are modular and can be extended in several directions. First, extension to multivariate ZITS is straightforward by widening the two output heads; the gate head would then model a joint sparsity pattern, potentially capturing cross-series zero-clustering dependencies. Second, conditioning mechanisms (e.g., class labels, calendar features, or covariate time series) can be integrated into both the encoder and the decoder to enable conditional ZITS generation, which is essential for retail demand simulation or IoT device modelling where device type or product category is a known covariate. Third, the temporal consistency loss can be generalised to integrate temporal components (such as trend and seasonality, similar to TimeVAE), higher-order autocorrelation lags, or cross-correlation between the gate and magnitude processes, further improving the realism of generated intermittency patterns. Finally, integration of differential privacy mechanisms would enable ZITS generation in settings where the underlying data contains sensitive personal information, such as individual device usage logs or patient-level healthcare records. While promising, evaluating these architectures on broader domains such as healthcare and insurance remains an important direction for future validation.


%%%%%%%%%%%%%%%%%%%%%%%%%%%%%%%%%%%%%%%%%%
\vspace{6pt} 

%%%%%%%%%%%%%%%%%%%%%%%%%%%%%%%%%%%%%%%%%%
%% optional
%\supplementary{The following supporting information can be downloaded at \linksupplementary{s1}, Figure S1: title; Table S1: title; Video S1: title.}




%%%%%%%%%%%%%%%%%%%%%%%%%%%%%%%%%%%%%%%%%%
\authorcontributions{Conceptualization, E.-R.A. and R.P.; methodology, E.-R.A. and R.P.; software, E.-R.A. and R.P.; validation, E.-R.A. and R.P.; formal analysis, E.-R.A. and R.P.; investigation, E.-R.A. and R.P.; data curation, E.-R.A. and R.P.; writing—original draft preparation, E.-R.A. and R.P.; writing—review and editing, E.-R.A. and R.P.; visualization, E.-R.A. and R.P.; supervision, E.-R.A. and R.P.; administration, E.-R.A. and R.P. All authors have read and agreed to the published version of the manuscript.}

\funding{This work is supported by the project "Romanian Hub for Artificial Intelligence-HRIA", Smart Growth, Digitization and Financial Instruments Program, MySMIS no. 334906.}

%% \institutionalreview{In this section, please add the Institutional Review Board Statement and approval number for studies involving humans or animals. You might choose to exclude this statement if the study did not require ethical approval. Please note that the Editorial Office might ask you for further information. Please add ``The study was conducted in accordance with the Declaration of Helsinki, and approved by the Institutional Review Board (or Ethics Committee) of NAME OF INSTITUTE (protocol code XXX and date of approval).” for studies involving humans OR ``The animal study protocol was approved by the Institutional Review Board (or Ethics Committee) of NAME OF INSTITUTE (protocol code XXX and date of approval).” for studies involving animals OR ``Ethical review and approval were waived for this study due to REASON (please provide a detailed justification).” OR ``Not applicable.” for studies not involving humans or animals.}

%% \informedconsent{Any research article describing a study involving humans should contain this statement. Please add ``Informed consent was obtained from all subjects involved in the study.'' OR ``Patient consent was waived due to REASON (please provide a detailed justification).'' OR ``Not applicable'' for studies not involving humans. You might also choose to exclude this statement if the study did not involve humans.

%%Written informed consent for publication must be obtained from participating patients who can be identified (including by the patients themselves). Please state ``Written informed consent has been obtained from the patient(s) to publish this paper.” if applicable.}

\dataavailability{
  The M5 forecasting dataset can be found at: \url{https://www.kaggle.com/competitions/m5-forecasting-accuracy/data}. 
  The code written in Python can be found at: \url{https://github.com/ArdeleanRichard/Zero-Inflated-Time-Series-Generation}
} 

%\durcstatement{Current research is limited to the [please insert a specific academic field, e.g., XXX], which is beneficial [share benefits and/or primary use] and does not pose a threat to public health or national security. Authors acknowledge the dual-use potential of the research involving xxx and confirm that all necessary precautions have been taken to prevent potential misuse. As an ethical responsibility, authors strictly adhere to relevant national and international laws about DURC. Authors advocate for responsible deployment, ethical considerations, regulatory compliance, and transparent reporting to mitigate misuse risks and foster beneficial outcomes.}

%% \acknowledgments{In this section, you can acknowledge any support given that is not covered by the author contribution or funding sections. This may include administrative and technical support, or donations in kind (e.g., materials used for experiments). Where GenAI has been used for purposes such as generating text, data, or graphics, or for study design, data collection, analysis, or interpretation of data, please add ``During the preparation of this manuscript/study, the author(s) used [tool name, version information] for the purposes of [description of use]. The authors have reviewed and edited the output and take full responsibility for the content of this publication.”}

\conflictsofinterest{The authors declare no conflicts of interest.} 


















%%%%%%%%%%%%%%%%%%%%%%%%%%%%%%%%%%%%%%%%%%
%% Optional
\appendixtitles{yes} % Leave argument "no" if all appendix headings stay EMPTY (then no dot is printed after "Appendix A"). If the appendix sections contain a heading then change the argument to "yes".
\appendixstart
\appendix


\section[\appendixname~\thesection]{Hyperparameter Optimisation}
\label{app:hyperparams}

Due to the computational cost of training generative models, a systematic but tractable hyperparameter search was conducted separately for the VAE and GAN variants of ZITS. For each model, a candidate set of configurations was defined over a discrete grid of hyperparameters (of approximately 1500 different configuration combinations). The VAE grid covered latent dimensionality, learning rate, and the loss weighting coefficients $\beta$, gate weight, reconstruction weight, and TC weight. The GAN grid covered the same structural parameters with the addition of the Adam $\beta_1$ momentum term and a feature-matching loss weight. 

Rather than exhaustively evaluating all combinations — which would amount to several hundred training runs per model — a random search strategy was employed, drawing 50 configurations per model from the full grid. Each configuration was trained for 80 epochs, after which 5,000 synthetic samples were generated and evaluated against the real data using the 21 distributional performance metrics. These metrics assessed fidelity across statistical moments (mean, standard deviation, skewness, kurtosis), quantile coverage, autocorrelation structure, and distributional divergence measures including Wasserstein distance, KL divergence, and MMD, as well as discriminative and predictive scores.

A well-known challenge in multi-metric evaluation is that no single configuration dominates across all criteria simultaneously — a run that minimises Wasserstein distance may rank poorly on autocorrelation fidelity, and vice versa. To resolve this, rank aggregation via the Borda count method was used to identify the configuration that performs most consistently well across the full set of metrics.

The Borda count is a positional voting method originating in social choice theory. Each metric is treated as a voter that casts a ranked ballot over all hyperparameter configurations, ordering them from best to worst according to that metric's value (lower being better in all cases). A configuration ranked $k-th$ out of $N$ candidates by a given metric receives a score of $N-k+1$ points from that voter. Final Borda scores are obtained by summing the points awarded across all metrics, and configurations are ranked by their total score. This procedure has the desirable property of rewarding consistent performance: a configuration that ranks in the top tier across many metrics will accumulate a high aggregate score, whereas a configuration that wins on one metric but performs poorly elsewhere will be penalised by the remaining voters. Tied metric values were assigned the minimum rank of the tied group prior to scoring. Configurations absent from a metric's ranking due to missing values were assigned a fractional score equal to the average of the unawarded positions, following the split-points convention.

The search also provides indirect evidence on the TC loss. Among the configurations shown in Tables~\ref{tab:hyperparams_iot_gan}--\ref{tab:hyperparams_m5_vae}, $\alpha_{tc}=0$ appears in 13 of 32 and $\alpha_{tc}>0$ in 19 of 32. In three of the four tables, the configuration with the lowest autocorrelation MAE has $\alpha_{tc}>0$, but it ranks lower overall under Borda aggregation.

The configuration with the highest Borda score was selected as the best hyperparameter setting for each model variant and dataset.

\begin{longtable}{l ccccc}
\caption{Results obtained by the ZITS-GAN model on the IoT data with multiple configurations
(columns are ordered from left to right based on Borda ranking -- only a subset is shown)}
\label{tab:hyperparams_iot_gan}\\

\toprule
\textbf{RunID} & \textbf{1} & \textbf{2} & \textbf{3} & \textbf{15} & \textbf{7} \\
\midrule
\endhead

\midrule
\multicolumn{6}{r}{\textit{Continued on next page}}\\
\endfoot

\bottomrule
\endlastfoot

$\mathit{latent\_dim}$   & 128    & 64     & 32     & 32     & 64     \\
$\mathit{lr}$            & 0.0002 & 0.0001 & 0.0001 & 0.0002 & 0.0002 \\
$\beta_0$      & 0.5    & 0.5    & 0.5    & 0.5    & 0      \\
$\alpha_{g}$  & 1      & 5      & 1      & 5      & 10     \\
$\alpha_{r}$  & 10     & 10     & 10     & 5      & 5      \\
$\alpha_{tc}$     & 0      & 0.5    & 2      & 1      & 0      \\
$\alpha_{fm}$     & 0.5    & 0.5    & 0.5    & 1      & 2      \\
$\mathit{train time (s)}$    & 322.8  & 322.5  & 346.1  & 348.6  & 427.4  \\

\midrule

$\Delta$ 0 ratio        & 0.024  & 0.034  & 0.063  & 0.011  & 0.002  \\
$\Delta$ mean            & 711.912 & 915.524 & 801.783 & 1207.302 & 398.853 \\
$\Delta$ std             & 910.615 & 1211.665 & 631.076 & 5474.926 & 290.094 \\
$\Delta$ skewness        & 0.540 & 0.348 & 0.597 & 0.499 & 2.895 \\
$\Delta$ Kurtosis        & 13.079 & 10.180 & 15.794 & 1.782 & 51.430 \\
$\Delta$Q10              & 0.000 & 0.000 & 0.000 & 0.000 & 0.000 \\
$\Delta$Q25              & 0.000 & 0.000 & 0.000 & 0.000 & 0.000 \\
$\Delta$Q50              & 0.000 & 0.000 & 0.000 & 0.000 & 0.000 \\
$\Delta$Q75              & 3568.698 & 3875.374 & 5040.737 & 1319.984 & 971.351 \\
$\Delta$Q90              & 2822.448 & 3620.289 & 2213.275 & 2867.933 & 1033.969 \\
$\Delta$Q95              & 1545.588 & 2344.104 & 356.607 & 18394.516 & 2267.978 \\
$\Delta$Q99              & 3424.546 & 1884.621 & 4819.143 & 25219.802 & 7661.044 \\
$\Delta$Qmean            & 1623.040 & 1674.913 & 1775.680 & 6828.891 & 1704.906 \\

\midrule

Autocorr. MAE           & 0.106 & 0.094 & 0.097 & 0.033 & 0.094 \\
Wasserstein dist.       & 2210.168 & 2290.182 & 1754.323 & 8842.176 & 2086.122 \\
KL divergence           & 0.155 & 0.143 & 0.152 & 1.008 & 0.130 \\
MMD                     & 0.002 & 0.002 & 0.002 & 0.002 & 0.002 \\
PS                      & 2376.839 & 2376.955 & 2376.864 & 2377.020 & 2376.931 \\
DS                      & 0.000 & 0.000 & 0.000 & 0.000 & 0.041 \\
LPS                     & 2376.940 & 2377.072 & 2376.950 & 2376.917 & 2376.919 \\
LDS                     & 0.250 & 0.194 & 0.211 & 0.005 & 0.259 \\

\end{longtable}



\begin{table}[htbp]
\centering
\caption{Results obtained by the ZITS-VAE model on the IoT data with multiple configurations (columns are ordered from left to right based on Borda ranking - only a subset is shown)}
\label{tab:hyperparams_iot_vae}
\scriptsize
\setlength{\tabcolsep}{1.5pt}
\begin{tabular*}{\textwidth}{@{\extracolsep{\fill}}l|cccccccccccccc}
\toprule
\textbf{RunID} &
\shortstack{\textbf{10}} &
\shortstack{\textbf{3}} &
\shortstack{\textbf{8}} &
\shortstack{\textbf{22}} &
\shortstack{\textbf{1}} &
\shortstack{\textbf{36}} &
\shortstack{\textbf{28}}  \\
\midrule
$\mathit{latent\_dim}$                         & 128      & 128       & 128       & 128      & 32        & 128      & 128              \\ \hline
$\mathit{lr}$                                  & 0.0005   & 0.0005    & 0.0005    & 0.001    & 0.0005    & 0.0001   & 0.0005           \\ \hline
$\beta_0$                            & 0.1      & 1         & 1         & 0.3      & 0.1       & 0.1      & 0.3              \\ \hline
$\alpha_{g}$                        & 10       & 5         & 5         & 10       & 10        & 5        & 10               \\ \hline
$\alpha_{r}$                        & 10       & 20        & 5         & 1        & 5         & 1        & 20               \\ \hline
$\alpha_{tc}$                           & 0.5      & 1         & 0         & 0        & 0         & 1        & 0                \\ \hline
$\mathit{train time (s)}$                         & 255.7    & 550.2     & 390.1     & 255.5    & 275.8     & 288.8    & 271.5            \\ \hline
$\Delta$ 0 ratio                    & 0.015    & 0.023     & 0.020     & 0.034    & 0.007     & 0.014    & 0.012            \\ \hline
$\Delta$ mean                       & 617.524  & 4178.317  & 3043.261  & 549.985  & 4516.269  & 170.717  & 1001.164         \\ \hline
$\Delta$ std                        & 795.339  & 8184.178  & 6332.825  & 503.951  & 8878.384  & 375.631  & 1507.860         \\ \hline
$\Delta$ skewness                   & 0.313    & 0.302     & 0.032     & 0.253    & 0.560     & 0.145    & 0.385            \\ \hline
$\Delta$ Kurtosis                   & 1.926    & 2.930     & 0.916     & 1.044    & 4.885     & 2.227    & 2.477            \\ \hline
$\Delta$Q10                         & 0.000    & 0.000     & 0.000     & 0.000    & 0.000     & 0.000    & 0.000            \\ \hline
$\Delta$Q25                         & 0.000    & 0.000     & 0.000     & 0.000    & 0.000     & 0.000    & 0.000            \\ \hline
$\Delta$Q50                         & 0.000    & 0.000     & 0.000     & 0.000    & 0.000     & 0.000    & 0.000            \\ \hline
$\Delta$Q75                         & 2830.178 & 7420.989  & 5425.259  & 3254.340 & 5584.768  & 1320.000 & 3566.197         \\ \hline
$\Delta$Q90                         & 2161.948 & 14578.944 & 9869.617  & 1424.636 & 18578.495 & 771.441  & 3645.603         \\ \hline
$\Delta$Q95                         & 2221.229 & 22659.868 & 16739.468 & 1195.924 & 26456.066 & 1445.158 & 4086.989         \\ \hline
$\Delta$Q99                         & 2245.594 & 39042.171 & 31947.538 & 1258.759 & 37046.733 & 2008.927 & 5097.523         \\ \hline
$\Delta$Qmean                       & 1351.279 & 11957.425 & 9140.269  & 1019.094 & 12523.723 & 792.218  & 2342.330         \\ \hline
\shortstack{Autocorr.\\MAE}         & 0.080    & 0.039     & 0.031     & 0.074    & 0.083     & 0.081    & 0.066            \\ \hline
\shortstack{Wasserstein\\distance}  & 1695.392 & 13970.431 & 10200.098 & 742.753  & 16727.510 & 775.787  & 3238.699         \\ \hline
\shortstack{KL\\divergence}         & 0.115    & 1.038     & 0.667     & 0.086    & 0.928     & 0.134    & 0.249            \\ \hline
MMD                                 & 0.002    & 0.002     & 0.002     & 0.002    & 0.002     & 0.002    & 0.002            \\ \hline
PS                                  & 2376.919 & 2376.904  & 2377.047  & 2376.890 & 2376.948  & 2376.989 & 2377.096         \\ \hline
DS                                  & 0.000    & 0.000     & 0.000     & 0.000    & 0.000     & 0.000    & 0.000            \\ \hline
LPS                                 & 2377.067 & 2376.934  & 2376.968  & 2376.911 & 2376.876  & 2376.910 & 2377.065         \\ \hline
LDS                                 & 0.060    & 0.099     & 0.077     & 0.096    & 0.099     & 0.185    & 0.059           \\

\bottomrule
\end{tabular*}
\end{table}







\begin{table}[htbp]
\centering
\caption{Results obtained by the ZITS-GAN model on the M5 data with multiple configurations (columns are ordered from left to right based on Borda ranking - only a subset is shown)}
\label{tab:hyperparams_m5_gan}
\scriptsize
\setlength{\tabcolsep}{1.5pt}
\begin{tabular*}{\textwidth}{@{\extracolsep{\fill}}l|cccccccccccccc}
\toprule
\textbf{RunID} &
\shortstack{\textbf{4}} &
\shortstack{\textbf{1}} &
\shortstack{\textbf{2}} &
\shortstack{\textbf{8}} &
\shortstack{\textbf{5}} &
\shortstack{\textbf{14}} &
\shortstack{\textbf{7}} &
\shortstack{\textbf{0}} &
\shortstack{\textbf{16}} \\
\midrule
$\mathit{latent\_dim}$                         & 32      & 128     & 128     & 64      & 32      & 64      & 64       & 64      & 64              \\ \hline
$\mathit{lr}$                                  & 0.0001  & 0.0001  & 0.0002  & 0.0001  & 0.0002  & 0.0002  & 5.00E-05 & 0.0001  & 0.0002          \\ \hline
$\beta_0$                            & 0       & 0       & 0       & 0.5     & 0       & 0       & 0.5      & 0       & 0               \\ \hline
$\alpha_{g}$                        & 1       & 10      & 10      & 10      & 5       & 10      & 10       & 1       & 1               \\ \hline
$\alpha_{r}$                        & 10      & 5       & 5       & 10      & 10      & 5       & 1        & 20      & 10              \\ \hline
$\alpha_{tc}$                           & 0       & 0.5     & 0       & 2       & 0.5     & 1       & 0        & 0       & 1               \\ \hline
$\alpha_{fm}$                           & 2       & 0.5     & 1       & 1       & 2       & 2       & 1        & 2       & 0.5             \\ \hline
$\mathit{train time (s)}$                         & 242.7   & 362.1   & 342.2   & 242.2   & 404.7   & 456.3   & 537.7    & 242.1   & 243.6           \\ \hline
$\Delta$ 0 ratio                    & 0.004   & 0.142   & 0.170   & 0.021   & 0.042   & 0.218   & 0.145    & 0.036   & 0.014           \\ \hline
$\Delta$ mean                       & 0.606   & 0.239   & 0.288   & 0.594   & 0.308   & 0.027   & 0.687    & 0.947   & 0.631           \\ \hline
$\Delta$ std                        & 2.601   & 1.462   & 1.495   & 2.602   & 0.561   & 0.744   & 2.427    & 0.969   & 2.642           \\ \hline
$\Delta$ skewness                   & 9.542   & 9.370   & 9.434   & 9.344   & 9.434   & 9.302   & 9.303    & 10.952  & 9.465           \\ \hline
$\Delta$ Kurtosis                   & 277.476 & 280.295 & 280.982 & 274.517 & 281.378 & 280.475 & 278.286  & 285.783 & 276.970         \\ \hline
$\Delta$Q10                         & 0.000   & 0.000   & 0.000   & 0.000   & 0.000   & 0.000   & 0.000    & 0.000   & 0.000           \\ \hline
$\Delta$Q25                         & 0.000   & 0.000   & 0.000   & 0.000   & 0.000   & 0.000   & 0.000    & 0.000   & 0.000           \\ \hline
$\Delta$Q50                         & 0.000   & 0.000   & 0.000   & 0.000   & 0.000   & 0.000   & 0.000    & 0.000   & 0.000           \\ \hline
$\Delta$Q75                         & 0.000   & 0.000   & 1.000   & 0.000   & 0.000   & 1.000   & 0.000    & 4.000   & 0.000           \\ \hline
$\Delta$Q90                         & 1.000   & 2.000   & 2.000   & 1.000   & 4.000   & 3.000   & 0.000    & 3.000   & 1.000           \\ \hline
$\Delta$Q95                         & 3.000   & 0.000   & 0.000   & 3.000   & 4.000   & 2.000   & 3.000    & 1.000   & 3.000           \\ \hline
$\Delta$Q99                         & 11.000  & 7.000   & 7.000   & 11.000  & 3.000   & 4.000   & 11.000   & 7.000   & 11.000          \\ \hline
$\Delta$Qmean                       & 2.143   & 1.286   & 1.429   & 2.143   & 1.571   & 1.429   & 2.000    & 2.143   & 2.143           \\ \hline
\shortstack{Autocorr.\\MAE}         & 0.063   & 0.171   & 0.127   & 0.060   & 0.123   & 0.022   & 0.059    & 0.077   & 0.046           \\ \hline
\shortstack{Wasserstein\\distance}  & 1.503   & 1.881   & 2.205   & 1.543   & 1.945   & 4.180   & 1.374    & 3.087   & 1.602           \\ \hline
\shortstack{KL\\divergence}         & 0.908   & 0.230   & 0.249   & 0.542   & 0.525   & 0.408   & 0.518    & 0.457   & 0.918           \\ \hline
MMD                                 & 0.002   & 0.002   & 0.002   & 0.002   & 0.002   & 0.002   & 0.002    & 0.002   & 0.002           \\ \hline
PS                                  & 1.372   & 1.316   & 1.333   & 1.349   & 1.380   & 1.329   & 1.392    & 1.372   & 1.410           \\ \hline
DS                                  & 0.000   & 0.004   & 0.000   & 0.000   & 0.000   & 0.000   & 0.146    & 0.026   & 0.000           \\ \hline
LPS                                 & 1.348   & 1.341   & 1.370   & 1.364   & 1.347   & 1.368   & 1.373    & 1.344   & 1.358           \\ \hline
LDS                                 & 0.297   & 0.458   & 0.447   & 0.324   & 0.456   & 0.487   & 0.398    & 0.469   & 0.233          \\

\bottomrule
\end{tabular*}
\end{table}






\begin{table}[htbp]
\centering
\caption{Results obtained by the ZITS-VAE model on the M5 data with multiple configurations (columns are ordered from left to right based on Borda ranking - only a subset is shown)}
\label{tab:hyperparams_m5_vae}
\scriptsize
\setlength{\tabcolsep}{1.5pt}
\begin{tabular*}{\textwidth}{@{\extracolsep{\fill}}l|cccccccccccccc}
\toprule
\textbf{RunID} &
\shortstack{\textbf{8}} &
\shortstack{\textbf{9}} &
\shortstack{\textbf{1}} &
\shortstack{\textbf{7}} &
\shortstack{\textbf{27}} &
\shortstack{\textbf{20}} &
\shortstack{\textbf{21}} &
\shortstack{\textbf{10}} &
\shortstack{\textbf{3}}  \\
\midrule
$\mathit{latent\_dim}$                         & 128     & 32      & 32      & 32      & 128     & 32      & 64      & 128     & 128             \\ \hline
$\mathit{lr}$                                  & 0.0005  & 0.0005  & 0.0005  & 0.0005  & 0.001   & 0.0001  & 0.0001  & 0.0005  & 0.0005          \\ \hline
$\beta_0$                            & 1       & 0.1     & 0.1     & 0.3     & 1       & 0.3     & 0.3     & 0.1     & 1               \\ \hline
$\alpha_{g}$                        & 5       & 5       & 10      & 10      & 10      & 10      & 5       & 10      & 5               \\ \hline
$\alpha_{r}$                        & 5       & 1       & 5       & 20      & 1       & 20      & 10      & 10      & 20              \\ \hline
$\alpha_{tc}$                           & 0       & 0.5     & 0       & 0.5     & 1       & 0       & 1       & 0.5     & 1               \\ \hline
$\mathit{train time (s)}$                         & 370.6   & 220.6   & 175.7   & 261.1   & 259.9   & 273.7   & 288.7   & 238.1   & 437.8           \\ \hline
$\Delta$ 0 ratio                    & 0.022   & 0.002   & 0.015   & 0.012   & 0.001   & 0.015   & 0.014   & 0.004   & 0.034           \\ \hline
$\Delta$ mean                       & 0.300   & 0.517   & 0.449   & 0.825   & 0.250   & 0.880   & 0.484   & 0.437   & 1.374           \\ \hline
$\Delta$ std                        & 1.296   & 0.942   & 1.139   & 0.434   & 1.952   & 0.231   & 0.931   & 2.339   & 0.407           \\ \hline
$\Delta$ skewness                   & 9.719   & 9.811   & 9.960   & 9.617   & 8.959   & 9.408   & 9.544   & 9.724   & 9.517           \\ \hline
$\Delta$ Kurtosis                   & 280.246 & 281.154 & 281.866 & 280.188 & 272.153 & 279.179 & 279.432 & 279.749 & 280.135         \\ \hline
$\Delta$Q10                         & 0.000   & 0.000   & 0.000   & 0.000   & 0.000   & 0.000   & 0.000   & 0.000   & 0.000           \\ \hline
$\Delta$Q25                         & 0.000   & 0.000   & 0.000   & 0.000   & 0.000   & 0.000   & 0.000   & 0.000   & 0.000           \\ \hline
$\Delta$Q50                         & 0.000   & 0.000   & 0.000   & 0.000   & 0.000   & 0.000   & 0.000   & 0.000   & 0.000           \\ \hline
$\Delta$Q75                         & 2.000   & 2.000   & 2.000   & 3.000   & 1.000   & 3.000   & 2.000   & 1.000   & 3.000           \\ \hline
$\Delta$Q90                         & 2.000   & 3.000   & 2.000   & 3.000   & 0.000   & 3.000   & 2.000   & 0.000   & 5.000           \\ \hline
$\Delta$Q95                         & 0.000   & 1.000   & 1.000   & 2.000   & 2.000   & 3.000   & 1.000   & 3.000   & 5.000           \\ \hline
$\Delta$Q99                         & 5.000   & 4.000   & 5.000   & 1.000   & 8.000   & 0.000   & 3.000   & 10.000  & 3.000           \\ \hline
$\Delta$Qmean                       & 1.286   & 1.429   & 1.429   & 1.286   & 1.571   & 1.286   & 1.143   & 2.000   & 2.286           \\ \hline
\shortstack{Autocorr.\\MAE}         & 0.046   & 0.061   & 0.064   & 0.058   & 0.023   & 0.067   & 0.068   & 0.024   & 0.058           \\ \hline
\shortstack{Wasserstein\\distance}  & 1.545   & 2.029   & 1.896   & 2.389   & 1.121   & 2.404   & 1.763   & 1.279   & 3.228           \\ \hline
\shortstack{KL\\divergence}         & 0.222   & 0.225   & 0.237   & 0.236   & 0.251   & 0.242   & 0.213   & 0.397   & 0.286           \\ \hline
MMD                                 & 0.002   & 0.002   & 0.002   & 0.002   & 0.002   & 0.002   & 0.002   & 0.002   & 0.002           \\ \hline
PS                                  & 1.370   & 1.326   & 1.372   & 1.373   & 1.373   & 1.295   & 1.344   & 1.338   & 1.345           \\ \hline
DS                                  & 0.000   & 0.008   & 0.000   & 0.000   & 0.000   & 0.000   & 0.000   & 0.000   & 0.142           \\ \hline
LPS                                 & 1.341   & 1.344   & 1.357   & 1.354   & 1.350   & 1.339   & 1.339   & 1.351   & 1.330           \\ \hline
LDS                                 & 0.396   & 0.416   & 0.407   & 0.424   & 0.275   & 0.404   & 0.388   & 0.286   & 0.397         \\

\bottomrule
\end{tabular*}
\end{table}



\begin{figure}[htbp]
  \centering
  \subfloat[ZITS-GAN.\label{fig:iot-zits-gan}]{%
    \includegraphics[width=0.48\textwidth]{fig/zits-gan-iot.png}%
  }%
  \hfill
  \subfloat[ZITS-VAE.\label{fig:iot-zits-vae}]{%
    \includegraphics[width=0.48\textwidth]{fig/zits-vae-iot.png}%
  }%
  \caption{Synthetically generated (orange) examples of the (a) ZITS-GAN and (b) ZITS-VAE proposed models for the IoT dataset.}
  \label{fig:iot-zits-side-by-side}
\end{figure}


\section[\appendixname~\thesection]{Run-Length Analysis}
\label{app:run-length}

To assess whether the proposed models reproduce the temporal clustering of zeros and non-zero observations, we compared run lengths in real and synthetic series. A run is a maximal contiguous sequence of time steps in the same state (zero or non-zero). We counted zero and non-zero runs separately across the series in each dataset and computed their empirical cumulative distribution functions (CDFs). Figures~\ref{fig:run-length-m5} and~\ref{fig:run-length-iot} show these CDFs for ZITS-GAN and ZITS-VAE; the horizontal axis is logarithmic to make both short and long runs visible.

For the M5 data, ZITS-GAN's CDFs closely track those of the real series for both zero and non-zero runs. ZITS-VAE also captures the broad zero-run distribution, but its non-zero-run CDF rises more quickly at short lengths, indicating a tendency to generate shorter non-zero runs than observed in the real data. The same overall contrast is visible on the IoT data: ZITS-GAN closely matches both run-length distributions, whereas ZITS-VAE shows a larger discrepancy. In particular, its zero-run CDF is lower at short lengths, consistent with longer zero runs, while its non-zero-run CDF is higher at short lengths, consistent with shorter non-zero runs. Thus, the run-length comparison supports the ability of ZITS-GAN to reproduce the temporal clustering in both datasets, while also showing that matching marginal zero ratios does not by itself ensure equally accurate zero and non-zero run lengths.

\begin{figure}[htbp]
  \centering
  \subfloat[ZITS-GAN.\label{fig:run-length-m5-gan}]{%
    \includegraphics[width=0.49\textwidth]{fig/run-length-m5-gan.png}%
  }
  \hfill
  \subfloat[ZITS-VAE.\label{fig:run-length-m5-vae}]{%
    \includegraphics[width=0.49\textwidth]{fig/run-length-iot-vae.png}%
  }
  \caption{Empirical CDFs of zero-run and non-zero-run lengths for real and synthetic M5 time series.}
  \label{fig:run-length-m5}
\end{figure}

\begin{figure}[htbp]
  \centering
  \subfloat[ZITS-GAN.\label{fig:run-length-iot-gan}]{%
    \includegraphics[width=0.49\textwidth]{fig/run-length-m5-gan.png}%
  }
  \hfill
  \subfloat[ZITS-VAE.\label{fig:run-length-iot-vae}]{%
    \includegraphics[width=0.49\textwidth]{fig/run-length-iot-vae.png}%
  }
  \caption{Empirical CDFs of zero-run and non-zero-run lengths for real and synthetic household IoT time series.}
  \label{fig:run-length-iot}
\end{figure}


\section[\appendixname~\thesection]{Ablation Study}
\label{app:ablation}

We conducted a controlled ablation study on the M5 dataset using the ZITS-VAE architecture. Five variants were evaluated: the full model; a model without the gate head; a model without the magnitude head; a model without the temporal consistency (TC) loss; and a model in which dilated convolution blocks in both the encoder and decoder were replaced by standard, non-dilated convolution blocks. For the single-head variants, the remaining head directly supplied the output. Each variant was trained in three independent runs with the same reduced computational budget (latent dimension 32, 25 epochs, batch size 128); the table reports arithmetic means across runs. The reported values exclude the real and synthetic metric components and include only comparison metrics.

The gate head has the clearest contribution to reproducing sparsity: the mean zero-ratio difference increases from 0.024 for the full model to 1.345 without the gate head, and the synthetic zero ratio in that variant is approximately zero. Removing the magnitude head also degrades value-distribution fidelity: the mean difference increases from 1.374 to 6.408, the Wasserstein distance from 3.372 to 21.087, and the KL divergence from 0.290 to 34.830. These results support the complementary roles of the two heads in this experiment.

The TC ablation shows little change in the M5 metrics: autocorrelation MAE is 0.060 with TC and 0.060 without it, while run-length Wasserstein distance is 0.550 and 0.549, respectively. Thus, this experiment does not show a measurable benefit from adding TC for this configuration. Similarly, replacing dilated convolutions with standard convolutions produces broadly comparable results: autocorrelation MAE is 0.060 for the full model and 0.060 without dilation, and zero-ratio differences are 0.024 and 0.024. This ablation does not establish that dilated convolutions outperform LSTMs; it compares them only with standard convolutions, and the results do not show a clear advantage over that alternative on M5.

\begin{center}
\scriptsize
\setlength{\tabcolsep}{3pt}
\begin{longtable}{@{}lrrrrr@{}}
\caption{Mean M5 ablation metrics across three runs. Lower values indicate smaller differences for the corresponding metrics.}\label{tab:m5-ablation}\\
\toprule
\textbf{Metric} & \textbf{Full} & \shortstack{\textbf{No gate}\\\textbf{head}} & \shortstack{\textbf{No magnitude}\\\textbf{head}} & \shortstack{\textbf{No TC}\\\textbf{loss}} & \shortstack{\textbf{No dilated}\\\textbf{convolutions}} \\
\midrule
\endfirsthead
\caption[]{Mean M5 ablation metrics across three runs (continued).}\\
\toprule
\textbf{Metric} & \textbf{Full} & \shortstack{\textbf{No gate}\\\textbf{head}} & \shortstack{\textbf{No magnitude}\\\textbf{head}} & \shortstack{\textbf{No TC}\\\textbf{loss}} & \shortstack{\textbf{No dilated}\\\textbf{convolutions}} \\
\midrule
\endhead
\midrule
\multicolumn{6}{r}{\textit{Continued on next page}}\\
\endfoot
\bottomrule
\endlastfoot
Zero-ratio difference & 0.024 & 1.345 & 0.086 & 0.023 & 0.024 \\
Mean difference & 1.374 & 0.609 & 6.408 & 1.383 & 1.381 \\
Standard deviation difference & 0.396 & 2.464 & 7.576 & 0.418 & 0.408 \\
Skewness difference & 9.586 & 8.769 & 10.739 & 9.586 & 9.586 \\
Kurtosis difference & 280.486 & 270.138 & 285.983 & 280.494 & 280.498 \\
Mean quantile difference & 2.286 & 2.238 & 10.143 & 2.381 & 2.286 \\
Q10 difference & 0.000 & 1.000 & 0.000 & 0.000 & 0.000 \\
Q25 difference & 0.000 & 1.000 & 0.000 & 0.000 & 0.000 \\
Q50 difference & 0.000 & 2.000 & 0.000 & 0.000 & 0.000 \\
Q75 difference & 3.000 & 1.000 & 23.000 & 3.000 & 3.000 \\
Q90 difference & 5.000 & 0.000 & 21.000 & 5.000 & 5.000 \\
Q95 difference & 5.000 & 2.000 & 18.000 & 5.333 & 5.000 \\
Q99 difference & 3.000 & 8.667 & 9.000 & 3.333 & 3.000 \\
Autocorrelation MAE & 0.060 & 0.082 & 0.051 & 0.060 & 0.060 \\
Wasserstein distance & 3.372 & 1.170 & 21.087 & 3.407 & 3.383 \\
KL divergence & 0.290 & 0.504 & 34.830 & 0.292 & 0.291 \\
MMD & 0.002 & 0.018 & 0.002 & 0.002 & 0.002 \\
Predictive score & 1.123 & 1.170 & 1.087 & 1.114 & 1.112 \\
Discriminative score & 0.036 & 0.095 & 0.171 & 0.000 & 0.030 \\
Long predictive score & 1.117 & 1.142 & 1.143 & 1.107 & 1.093 \\
Long discriminative score & 0.447 & 0.500 & 0.495 & 0.445 & 0.452 \\
Run-length Wasserstein distance & 0.550 & 152.195 & 0.867 & 0.549 & 0.564 \\
\end{longtable}
\end{center}


%%%%%%%%%%%%%%%%%%%%%%%%%%%%%%%%%%%%%%%%%%
%\isPreprints{}{% This command is only used for ``preprints''.
\begin{adjustwidth}{-\extralength}{0cm}
%} % If the paper is ``preprints'', please uncomment this parenthesis.
%\printendnotes[custom] % Un-comment to print a list of endnotes

\reftitle{References}

% References must be numbered in order of appearance in the text (including citations in tables and legends) and listed individually at the end of the manuscript. We recommend preparing the references with a bibliography software package, such as EndNote, ReferenceManager or Zotero, to avoid typing mistakes and duplicated references. Include the digital object identifier (DOI) for all references where available.
% Please provide either the correct journal abbreviation (e.g. according to the “The list of Title Word Abbreviations” https://portal.issn.org/ltwa) or the full name of the journal. 
% Citations and references in the Supplementary Materials are permitted provided that they also appear in the reference list here. 
% In the text, reference numbers should be placed in square brackets [ ] and placed before the punctuation; for example, [1], [1–3] or [1,3]. For embedded citations in the text with pagination, use both parentheses and brackets to indicate the reference number and page numbers; for example, [5] (p. 10), or [6] (pp. 101–105).

%=====================================
% References, variant A: external bibliography
%=====================================
\nocite{*}
\externalbibliography{yes}
\bibliography{bibliography}

%%%%%%%%%%%%%%%%%%%%%%%%%%%%%%%%%%%%%%%%%%
%% for journal Sci
%\reviewreports{\\
%Reviewer 1 comments and authors’ response\\
%Reviewer 2 comments and authors’ response\\
%Reviewer 3 comments and authors’ response
%}
%%%%%%%%%%%%%%%%%%%%%%%%%%%%%%%%%%%%%%%%%%
\PublishersNote{}
%\isPreprints{}{% This command is only used for ``preprints''.
\end{adjustwidth}
%} % If the paper is ``preprints'', please uncomment this parenthesis.
\end{document}

